\documentclass[10pt]{article}

\usepackage[T1]{fontenc}
\usepackage[utf8]{inputenc}
\usepackage{lmodern}
\usepackage{microtype}
\usepackage[a4paper,margin=1in]{geometry}

\usepackage{amsmath,amssymb,amsthm,mathtools}
\usepackage{bm}
\usepackage{booktabs}
\usepackage{array}
\usepackage{tabularx}
\usepackage{longtable}
\usepackage{enumitem}
\usepackage{xcolor}
\usepackage{graphicx}
\usepackage{url}
\usepackage{natbib}
\usepackage{hyperref}
\usepackage[nameinlink,capitalise,noabbrev]{cleveref}

\definecolor{deepblue}{HTML}{173B57}
\definecolor{softblue}{HTML}{EAF3F8}
\definecolor{deepgreen}{HTML}{235D3A}
\definecolor{softgreen}{HTML}{EAF5EE}
\definecolor{deeporange}{HTML}{9C4E00}
\definecolor{softorange}{HTML}{FFF1DF}
\definecolor{draftred}{HTML}{8F1D1D}
\definecolor{draftbg}{HTML}{FFF4F4}
\definecolor{openpurple}{HTML}{5B3F8C}
\definecolor{softpurple}{HTML}{F2ECFA}
\definecolor{graytext}{HTML}{4B5563}

\hypersetup{
  colorlinks=true,
  linkcolor=deepblue,
  citecolor=deepgreen,
  urlcolor=deepblue,
  pdftitle={Nuclear-Norm Best Approximation at Rank-Deficient Matrix-Line Crossings},
  pdfauthor={Anonymous manuscript}
}

\setlength{\parindent}{0pt}
\setlength{\parskip}{0.55em}
\setlist[itemize]{topsep=0.35em,itemsep=0.2em}
\setlist[enumerate]{topsep=0.35em,itemsep=0.25em}
\renewcommand{\arraystretch}{1.15}

\newtheoremstyle{reader}
  {0.8em}{0.8em}
  {\itshape}{}
  {\bfseries\color{deepblue}}{.}
  {0.5em}{}
\theoremstyle{reader}
\newtheorem{theorem}{Theorem}[section]
\newtheorem{proposition}[theorem]{Proposition}
\newtheorem{lemma}[theorem]{Lemma}
\newtheorem{corollary}[theorem]{Corollary}

\theoremstyle{definition}
\newtheorem{definition}[theorem]{Definition}
\newtheorem{assumption}[theorem]{Assumption}
\newtheorem{example}[theorem]{Example}

\theoremstyle{remark}
\newtheorem{remark}[theorem]{Remark}

\newenvironment{proofidea}
  {\par\smallskip\begin{center}\begin{minipage}{0.94\textwidth}
   {\color{deepblue}\hrule}\vspace{0.35em}
   \textbf{\color{deepblue}Proof idea.}\ \color{black}\ignorespaces}
  {\par\vspace{0.35em}{\color{deepblue}\hrule}
   \end{minipage}\end{center}\smallskip}

\newenvironment{intuition}
  {\par\smallskip\begin{center}\begin{minipage}{0.94\textwidth}
   {\color{deepgreen}\hrule}\vspace{0.35em}
   \textbf{\color{deepgreen}What this result does.}\ \color{black}\ignorespaces}
  {\par\vspace{0.35em}{\color{deepgreen}\hrule}
   \end{minipage}\end{center}\smallskip}

\newenvironment{boundary}
  {\par\smallskip\begin{center}\begin{minipage}{0.94\textwidth}
   {\color{deeporange}\hrule}\vspace{0.35em}
   \textbf{\color{deeporange}Boundary of the claim.}\ \color{black}\ignorespaces}
  {\par\vspace{0.35em}{\color{deeporange}\hrule}
   \end{minipage}\end{center}\smallskip}

\newenvironment{openproblem}
  {\par\smallskip\begin{center}\begin{minipage}{0.94\textwidth}
   {\color{openpurple}\hrule}\vspace{0.35em}
   \textbf{\color{openpurple}Open problem.}\ \color{black}\ignorespaces}
  {\par\vspace{0.35em}{\color{openpurple}\hrule}
   \end{minipage}\end{center}\smallskip}

\newcommand{\proved}{\textcolor{deepgreen}{\textbf{PROVED}}}
\newcommand{\localproved}{\textcolor{deeporange}{\textbf{PROVED UNDER LOCAL ASSUMPTIONS}}}
\newcommand{\observed}{\textcolor{deepblue}{\textbf{NUMERICALLY OBSERVED}}}
\newcommand{\openstatus}{\textcolor{openpurple}{\textbf{OPEN}}}

\newcommand{\ip}[2]{\left\langle #1,#2\right\rangle}
\newcommand{\R}{\mathbb{R}}
\newcommand{\norm}[1]{\left\lVert #1\right\rVert}
\newcommand{\snorm}[1]{\left\lVert #1\right\rVert_{2}}
\newcommand{\fnorm}[1]{\left\lVert #1\right\rVert_{\mathrm F}}
\newcommand{\nnorm}[1]{\left\lVert #1\right\rVert_{*}}
\newcommand{\argmin}{\operatorname*{arg\,min}}
\newcommand{\argmax}{\operatorname*{arg\,max}}
\newcommand{\diag}{\operatorname{diag}}
\newcommand{\rank}{\operatorname{rank}}
\newcommand{\range}{\operatorname{range}}
\newcommand{\Retr}{\operatorname{Retr}}
\newcommand{\Gap}{\operatorname{Gap}}
\newcommand{\Sc}{S_{\mathrm c}}
\newcommand{\Kcan}{K_{\mathrm{can}}}
\newcommand{\Pstar}{\mathcal P^\star}
\newcommand{\Ctheta}{\mathcal C_{\Theta}}

\DeclareMathOperator{\tr}{tr}
\DeclareMathOperator{\dist}{dist}
\DeclareMathOperator{\sign}{sign}

\crefname{theorem}{theorem}{theorems}
\crefname{proposition}{proposition}{propositions}
\crefname{lemma}{lemma}{lemmas}
\crefname{corollary}{corollary}{corollaries}
\crefname{definition}{definition}{definitions}
\crefname{assumption}{assumption}{assumptions}
\crefname{example}{example}{examples}
\crefname{remark}{remark}{remarks}

\hypersetup{pdftitle={Nuclear-Norm Best Approximation at Rank-Deficient Matrix-Line Crossings}}
\setlength{\parskip}{0.35em}
\setlist[itemize]{topsep=0.25em,itemsep=0.1em}
\setlist[enumerate]{topsep=0.25em,itemsep=0.15em}

\title{\textbf{Nuclear-Norm Best Approximation at Rank-Deficient Matrix-Line Crossings}\\
\large Certificate Geometry and Singular Smoothing}
\author{Anonymous manuscript}
\date{August 4, 2026}

\begin{document}
\maketitle

\begin{abstract}
We study nuclear-norm best approximation from a matrix line and its dual
spectral-norm ball intersected with one hyperplane.  At a rank-deficient
minimizer, exact certificates form a tangency slice of the nuclear-norm
subdifferential.  We classify its feasibility and uniqueness, derive its
minimum-Frobenius member by singular-value water-filling, and characterize
when clipping is unnecessary.  An admissible Fenchel family unifies two
centers: the Huber regularizer has the water-filled center, whereas
pseudo-Huber generally selects a different smooth center.  A polynomial
derivative tail of order \(p\) yields boundary exponent \(p/(p+1)\), while
finite saturation yields a linear law.  For arbitrary matrix lines, we derive
the regular expansion and analytic multiplier and direction expansions at
strict rank-deficient minimizers of arbitrary corank, with computable first
corrections.  Compact-polar stationarity produces cubic/quadratic
superconvergence and an explicit quintic/quartic refinement after a second
cancellation.  At transverse square corank-one boundaries, pseudo-Huber obeys
the exact two-thirds law, and a moving-crossing normal form gives a uniform
cubic balance through the strict-to-boundary transition.  Finally, the
compact-polar selection fails on a nonempty open set even when the approximant
and exact certificate are unique.
\end{abstract}

\noindent\textbf{Keywords.}
Nuclear norm; best approximation; spectral norm; Birkhoff--James
orthogonality; polar factor; Fenchel smoothing; singular perturbation.

\noindent\textbf{MSC 2020.}
15A18, 15A60, 41A50, 41A52, 47A30, 49J52, 65F35, 65K10.

\section{Introduction}
\label{sec:intro}

Let \(G,\Theta\in\mathbb R^{m\times n}\), with \(\Theta\ne0\).  Our primary
object is the nuclear-norm distance from \(G\) to the one-dimensional matrix
subspace generated by \(\Theta\):
\begin{equation}
 d^\star
 =\min_{\lambda\in\mathbb R}\phi(\lambda),
 \qquad
 \phi(\lambda)=\|G+\lambda\Theta\|_*
 =\operatorname{dist}_{\|\cdot\|_*}
   \bigl(G,\operatorname{span}\{\Theta\}\bigr).
 \label{eq:D}
\end{equation}
Classical distance duality identifies the same value with
\begin{equation}
 p^\star=
 \max_{\Phi}
 \left\{\langle G,\Phi\rangle:
 \|\Phi\|_2\le1,\ \langle\Theta,\Phi\rangle=0\right\}.
 \label{eq:P}
\end{equation}
Thus \(d^\star=p^\star\): the dual unit ball is the spectral-norm ball, and
the annihilator of \(\operatorname{span}\{\Theta\}\) is one hyperplane.  We
study what happens when the affine matrix line
\(\lambda\mapsto G+\lambda\Theta\) meets the rank-deficient locus.  The
distinctive difficulty is that the natural point derivative is expressed by
the compact polar factor.  It is smooth on fixed-rank strata, but under
general rank-changing perturbations the zero-kernel completion need not have a
continuous extension; see, e.g., the classical polar-decomposition framework
of \citet{higham1986polar}.

When \(\Theta=uv^\top\), \eqref{eq:P} is the tangent LMO on the smooth
spectral sphere.  Norm-ball LMOs underlie matrix-aware first-order methods
\citep{pethick2025lmo}.  \citet{xie2026controlled} introduced the corresponding
compact-polar bisection, and \citet[Appendix P.5]{li2026intrinsic} interpreted
it as the spectral-sphere instance of an intrinsic LMO.  Neither work treats
rank-deficient crossings.

\paragraph{Classical boundary and the new object.}
Optimality in \eqref{eq:D} is classical Birkhoff--James orthogonality.
Distance duality, matrix-norm subdifferentials, and the associated criteria
are established in
\citep{singer1970best,watson1992subdifferential,zietak1993faces,
watson1993kyfan,bhatia1999orthogonality,grover2014distance,
johnston2022trace,grover2026kyfan}; we claim no novelty for those facts.
The new object is the tangency slice
\(\partial\|A^\star\|_*\cap\Theta^\perp\), its spectral centers, and the
singular asymptotics of their regularization paths.

Canonical approximant selection by Schatten limits is also classical:
\citet{legg1985canonical} treat \(p\downarrow1\), while
\citet{zietak2017strict} surveys the \(p\uparrow\infty\) strict-spectral
counterpart, including its conjectural general form.  Our selection instead
occurs on the equality-constrained certificate face.
If \(A=U_r\Sigma_rV_r^\top\) is a compact singular value decomposition,
the compact polar factor is
\begin{equation}
 \Sc(A)=U_rV_r^\top.
 \label{eq:compact-polar}
\end{equation}
At full rectangular rank, \(\phi\) is differentiable and
\(\phi'(\lambda)=\langle\Theta,\Sc(G+\lambda\Theta)\rangle\).  This suggests
the ordinary scalar equation
\begin{equation}
 h_{\rm c}(\lambda):=
 \langle\Theta,\Sc(G+\lambda\Theta)\rangle=0.
 \label{eq:compact-root}
\end{equation}
At full rectangular rank this equation is valid.  When the affine matrix loses
rank, it need not remain valid, because the complete subdifferential is
\begin{equation}
 \partial\|A\|_*
 =\left\{
 U_rV_r^\top+U_0KV_0^\top:\ \|K\|_2\le1
 \right\},
 \label{eq:full-subdiff}
\end{equation}
where \(U_0,V_0\) span the left and right null spaces.  The compact factor is
only the choice \(K=0\).  The omitted kernel block can be exactly what is
needed to enforce \(\langle\Theta,\Phi\rangle=0\).  Consequently,
\(h_{\rm c}\) can be monotone and still jump across zero.

Thus the LMO and nuclear-norm duality remain exact; only the replacement of a
set-valued optimality graph by one compact selection can fail.

\subsection{Contributions and novelty boundary}

The contributions, in the order used by the analysis, are as follows.

\begin{enumerate}
 \item Starting from the classical interval
 \(\partial\phi=[a-\beta,a+\beta]\)
 \citep{watson1992subdifferential,lewis1995convex}, we classify feasibility
 and uniqueness of the tangency slice.  Its minimum-Frobenius member follows
 by singular-value water-filling; clipping is universally unnecessary exactly
 for scalar multiples of partial isometries.

 \item An admissible Fenchel family unifies the hard and smooth centers:
 Huber gives \(K_F\), pseudo-Huber gives \(K_R\), and we characterize their
 agreement.  A polynomial derivative tail of order \(p>1\) yields exponent
 \(p/(p+1)\), while finite saturation gives a linear law.  The novelty beyond
 standard smoothing principles
 \citep{nesterov2005smooth,lobos2025smooth,attouch1984variational} is the
 constrained spectral center and its coupling to the singular multiplier.

 \item We prove the regular expansion and, at every strict rank-deficient
 minimizer, analytic multiplier and returned-direction expansions without
 shape, normal-rank, or corank restrictions.  Compact-root instances have
 computable cubic/quadratic superconvergence and a quintic/quartic refinement.
 At transverse square corank-one boundaries, pseudo-Huber has the exact
 two-thirds law and a uniform cubic strict-to-boundary crossover.  Our new
 content is the matrix-line exponents, constants, and crossover, rather than
 the existence of fractional central-path rates \citep{basu2022central}.

 \item Finally, a strict \(2\times2\) witness extends to a nonempty open
 compact-polar failure set in every square dimension at least two, even with
 unique approximant and certificate.  This motivates retaining the full
 slice.  Rank-one normals also admit an SVD-free crossing formula.
\end{enumerate}

An open failure set implies neither an event-frequency law nor an outer
algorithmic or learning-performance gain; those require additional models.

\subsection{Organization}

\Cref{sec:duality,sec:failure} give the scalar interval, structural failure,
and certificate geometry; \cref{sec:smoothing,sec:rates} develop the Fenchel
centers and local laws; \cref{sec:numerics} verifies the constants.  Auxiliary
level-set geometry and certificate repair are collected in the supplement.

\section{Exact duality and scalar optimality}
\label{sec:duality}

We use the Frobenius pairing
\(\langle X,Y\rangle=\operatorname{tr}(X^\top Y)\).  The spectral and
nuclear norms are dual.  For a compact SVD \(A=U_r\Sigma_rV_r^\top\), let
\(U_0,V_0\) be orthonormal null-space bases; a zero-dimensional factor is
allowed.  The standard formula \eqref{eq:full-subdiff} will be used without a
choice of bases inside repeated positive singular spaces.

\begin{theorem}[Strong duality and the complete solution face]
\label{thm:T1}
For arbitrary \(G\in\mathbb R^{m\times n}\) and nonzero
\(\Theta\in\mathbb R^{m\times n}\),
\begin{equation}
 p^\star
 =\max_{\substack{\|\Phi\|_2\le1\\
                   \langle\Theta,\Phi\rangle=0}}
   \langle G,\Phi\rangle
 =\min_{\lambda\in\mathbb R}\|G+\lambda\Theta\|_*.
 \label{eq:strong-duality}
\end{equation}
The dual minimizer set is nonempty and compact.  If \(\lambda^\star\) is any
dual minimizer and \(A^\star=G+\lambda^\star\Theta\), then all primal
solutions are
\begin{equation}
 \mathcal P^\star
 =\partial\|A^\star\|_*\cap\Theta^\perp.
 \label{eq:solution-face}
\end{equation}
Every dual minimizer satisfies
\begin{equation}
 |\lambda^\star|
 \le \frac{2\|G\|_*}{\|\Theta\|_*}.
 \label{eq:lambda-bound}
\end{equation}
\end{theorem}

\begin{proof}
The support function of the spectral unit ball is the nuclear norm.  Since
\(\Phi=0\) is feasible and lies in the interior of the ball,
finite-dimensional Slater--Fenchel duality gives
\eqref{eq:strong-duality}.  The reverse triangle inequality
\(\|G+\lambda\Theta\|_*\ge
|\lambda|\|\Theta\|_*-\|G\|_*\) gives coercivity and, by comparison with
\(\lambda=0\), \eqref{eq:lambda-bound}.  Finally, equality in the support
function is equivalent to
\(\Phi\in\partial\|A^\star\|_*\); imposing tangency yields exactly
\eqref{eq:solution-face}.
\end{proof}

At a fixed \(\lambda\), set
\begin{equation}
 A=G+\lambda\Theta=U_r\Sigma_rV_r^\top,
 \qquad
 a=\langle\Theta,U_rV_r^\top\rangle,
 \qquad
 C=U_0^\top\Theta V_0,
 \qquad
 \beta=\|C\|_*.
 \label{eq:a-C-beta-main}
\end{equation}

\begin{theorem}[Exact scalar interval and compact-root criterion]
\label{thm:interval}
For \(\phi(\lambda)=\|G+\lambda\Theta\|_*\),
\begin{equation}
 \partial\phi(\lambda)=[a-\beta,a+\beta].
 \label{eq:scalar-interval}
\end{equation}
Consequently, at this specified \(\lambda\):
\begin{enumerate}
 \item \(\lambda\) is dual optimal if and only if \(|a|\le\beta\);
 \item the compact equation \eqref{eq:compact-root} holds if and only if
 \(a=0\);
 \item dual optimality holds but compact stationarity fails if and only if
 \(0<|a|\le\beta\).
\end{enumerate}
Globally, \eqref{eq:compact-root} has no solution if and only if
\(a(\lambda)\ne0\) for every \(\lambda\in\argmin\phi\).  If this happens,
the dual minimizer is unique and its affine matrix is rank deficient.
\end{theorem}

\begin{proof}
The affine chain rule and \eqref{eq:full-subdiff} give
\(\partial\phi(\lambda)=\{a+\langle C,K\rangle:\|K\|_2\le1\}\), whose
endpoints are \(a\pm\|C\|_*\).  This proves
\eqref{eq:scalar-interval} and the three pointwise claims.  If two minimizers
existed, convexity would make \(\phi\) constant between them; at every
interior point \(\partial\phi=\{0\}\), hence \(a=\beta=0\) and a compact root
would exist.  Thus no-root failure forces uniqueness.  Full rectangular rank
would give \(\beta=0\), and optimality would again force \(a=0\), so the
unique no-root minimizer is rank deficient.
\end{proof}

\begin{remark}[Relative-interior nondegeneracy]
\label{rem:strict-complementarity-main}
At an optimum, \(|a|<\beta\) is exactly
\(0\in\operatorname{ri}\partial\phi(\lambda^\star)\); equivalently, the
tangency hyperplane meets the relative interior of the kernel-block spectral
unit ball.  At \(|a|=\beta\) it meets only the exposed endpoint face
\eqref{eq:endpoint-face}.  This parallels convergence and limiting-center
phenomena for degenerate semidefinite central paths
\citep{halicka2002convergence,halicka2005limiting}, while fractional algebraic
rates are documented by \citet{basu2022central}.  Our rates follow directly
from the matrix-line normal forms in \cref{sec:rates}.
\end{remark}

The compact selection is nevertheless nondecreasing.  Indeed, for
\(\lambda_1<\lambda_2\), each compact polar factor belongs to the corresponding
nuclear-norm subdifferential, and monotonicity of the convex subdifferential
graph gives \(h_{\rm c}(\lambda_1)\le h_{\rm c}(\lambda_2)\).  The obstruction
is a jump, not a reversal of monotonicity.

\begin{remark}[Exact interval bisection]
With \(L=a-\beta\) and \(U=a+\beta\), \(U<0\) moves the bracket to the
right, \(L>0\) moves it to the left, and \(L\le0\le U\) certifies a
minimizer.  The initial bracket is \eqref{eq:lambda-bound}; on termination one
applies the completion \eqref{eq:KF-main}.  Unlike compact-polar bisection,
this test cannot jump over a minimizer, but it requires a numerical-rank
decision.  The elementary convergence proof is recorded in the supplement.
\end{remark}

\section{Structural failure and certificate geometry}
\label{sec:failure}

The rank-one normal from the motivating tangent oracle supplies a transparent
witness.  The structural result below then perturbs that witness in the full
space of matrix normals; the completion theory is rectangular and makes no
rank assumption on \(\Theta\).

\subsection{A witness and full-space openness}

Let \(b_1,b_2,c>0\), \(d\in\mathbb R\), and
\begin{equation}
 \Theta=e_1e_1^\top,
 \qquad
 G=\begin{pmatrix}d&b_1\\b_2&c\end{pmatrix}.
 \label{eq:2by2-data}
\end{equation}
Writing \(x=d+\lambda\), set
\(A(x)=\bigl(\begin{smallmatrix}x&b_1\\b_2&c\end{smallmatrix}\bigr)\).

\begin{theorem}[A strict \(2\times2\) family with no compact root]
\label{thm:open-2by2}
If \(0<b_1b_2\le c^2\), then
\begin{equation}
 x^\star=\frac{b_1b_2}{c}
 \label{eq:2by2-minimizer}
\end{equation}
is the unique minimizer of \(\|A(x)\|_*\), and the compact scalar selection
does not attain zero.  At the singular minimizer,
\begin{equation}
 \left\langle e_1e_1^\top,\Sc(A(x^\star))\right\rangle
 =\frac{b_1b_2}{\sqrt{(b_1^2+c^2)(b_2^2+c^2)}}>0.
 \label{eq:2by2-jump}
\end{equation}
The strict condition \(0<b_1b_2<c^2\) is open in the positive parameter orthant.
\end{theorem}

\begin{proof}
For a \(2\times2\) matrix,
\(\|A\|_*^2=\|A\|_F^2+2|\det A|\).  Since
\(\det A(x)=xc-b_1b_2\),
\begin{equation}
 \|A(x)\|_*^2=
 \begin{cases}
 (x-c)^2+(b_1+b_2)^2,&x<x^\star,\\
 (x+c)^2+(b_1-b_2)^2,&x>x^\star.
 \end{cases}
 \label{eq:2by2-branches}
\end{equation}
Below the crossing, the derivative of the nuclear norm is
\((x-c)/\sqrt{(x-c)^2+(b_1+b_2)^2}\).  Above it, the derivative is
\((x+c)/\sqrt{(x+c)^2+(b_1-b_2)^2}\).  They are respectively negative and
positive because \(x^\star\le c\), proving uniqueness.
Moreover
\[
 A(x^\star)=
 \begin{pmatrix}b_1\\c\end{pmatrix}
 \begin{pmatrix}b_2/c&1\end{pmatrix},
\]
whose compact polar has the entry in \eqref{eq:2by2-jump}.  The compact
scalar therefore jumps from a negative to a positive value without attaining
zero.
\end{proof}

The equality case \(b_1b_2=c^2\) will be the boundary-kink model in
\cref{sec:rates}.  Strict instances give a stronger statement than relative
openness inside the rank-one manifold.  Write
\[
 \mathbb S_F^{n^2-1}
 =\{\Theta\in\mathbb R^{n\times n}:\|\Theta\|_F=1\}.
\]

\begin{theorem}[Full-space openness of strict compact-polar failure]
\label{thm:open-all-n}
For every \(n\ge2\), there are a nonempty relatively open set
\[
 \mathcal O_n\subset \mathrm{GL}_n\times\mathbb S_F^{n^2-1}
\]
and a real-analytic map
\((G,\Theta)\mapsto\lambda_{\rm s}(G,\Theta)\) such that
\(A_{\rm s}=G+\lambda_{\rm s}\Theta\) has rank \(n-1\).  With unit left and
right null vectors \(u_{\rm s},v_{\rm s}\), define
\begin{equation}
 a_{\rm s}=\langle\Theta,\Sc(A_{\rm s})\rangle,
 \qquad
 \beta_{\rm s}=|u_{\rm s}^\top\Theta v_{\rm s}|.
 \label{eq:full-open-ab}
\end{equation}
Then
\begin{equation}
 0<|a_{\rm s}|<\beta_{\rm s}.
 \label{eq:strict-failure}
\end{equation}
The crossing \(\lambda_{\rm s}\) is the unique global minimizer and satisfies
the sharp-growth bound
\begin{equation}
 \|G+\lambda\Theta\|_*
 \ge \|A_{\rm s}\|_*
 +\bigl(\beta_{\rm s}-|a_{\rm s}|\bigr)
  |\lambda-\lambda_{\rm s}|
 \quad(\lambda\in\mathbb R).
 \label{eq:sharp-growth}
\end{equation}
The compact equation \eqref{eq:compact-root} has no solution.  In particular,
strict failure has nonempty open interior in the full normalized data space,
which includes full-rank matrix normals.
\end{theorem}

\begin{proof}
For \(n=2\), take
\[
 G_0=\begin{pmatrix}0&1\\1&2\end{pmatrix},
 \qquad \Theta_0=e_1e_1^\top,
 \qquad \lambda_0=\frac12.
\]
Then \(G_0\) is invertible and
\(A_0=\tfrac12(1,2)^\top(1,2)\).  Its normalized range and null vectors are
\(w=(1,2)^\top/\sqrt5\) and \(z=(-2,1)^\top/\sqrt5\), respectively, so
\[
 a_0=\langle\Theta_0,ww^\top\rangle=\frac15,
 \qquad \beta_0=|z^\top\Theta_0z|=\frac45.
\]
For \(n>2\), embed \(G_0\) as the leading block of
\(\operatorname{diag}(G_0,d_3,\ldots,d_n)\), with \(d_j>0\), and extend
\(\Theta_0\) by zeros.  The crossing has corank one and retains these strict
values.

Set \(F(G,\Theta,\lambda)=\det(G+\lambda\Theta)\).  At the witness,
\(\partial_\lambda F\ne0\), because at a corank-one matrix the adjugate is a
nonzero multiple of \(v_{\rm s}u_{\rm s}^\top\) and
\(u_{\rm s}^\top\Theta v_{\rm s}\ne0\).  The analytic implicit-function
theorem, restricted to the Frobenius sphere, therefore continues the crossing
as \(\lambda_{\rm s}(G,\Theta)\).  After shrinking the neighborhood,
\(\partial_\lambda F\ne0\).  Jacobi's formula gives
\[
 \partial_\lambda\det A=\tr\!\left(\operatorname{adj}(A)\Theta\right),
\]
whereas \(\operatorname{adj}(A)=0\) whenever \(\rank A\le n-2\).
Thus a zero of \(F\) with \(\partial_\lambda F\ne0\) has rank exactly
\(n-1\), so the continued crossing remains corank one.  The
partial polar factor and the null projectors are continuous along this
fixed-rank crossing manifold, so \(a_{\rm s}\) and \(\beta_{\rm s}\) are
continuous and the strict inequalities persist.

At the crossing, \cref{thm:interval} gives
\(\partial\phi(\lambda_{\rm s})=[a_{\rm s}-\beta_{\rm s},
a_{\rm s}+\beta_{\rm s}]\).  Apply the subgradient inequality with the right
endpoint for \(\lambda>\lambda_{\rm s}\) and the left endpoint for
\(\lambda<\lambda_{\rm s}\); both imply \eqref{eq:sharp-growth}.  Thus the
minimizer is globally unique.  A compact-polar root anywhere would provide a
zero scalar subgradient and hence another minimizer; at \(\lambda_{\rm s}\)
the compact scalar equals \(a_{\rm s}\ne0\).  No root exists.  Finally,
full-rank normals are dense in \(\mathbb S_F^{n^2-1}\), so \(\mathcal O_n\)
contains such normals.
\end{proof}

The implicit-function theorem asserts local uniqueness of the continued
\emph{crossing}, not global uniqueness of all determinant roots; small
general-normal perturbations may create additional far-away crossings.
Equation \eqref{eq:sharp-growth} is what makes the continued crossing the
unique global \emph{minimizer}.  The theorem is topological and makes no
frequency claim.

For the rank-one normal appearing in the tangent-oracle application, the
crossing and its half-width admit an SVD-free formula.  Let
\begin{equation}
 \Theta=uv^\top,
 \qquad \|u\|_2=\|v\|_2=1.
 \label{eq:rank-one-normal}
\end{equation}

\begin{proposition}[Rank-one crossing formula]
\label{prop:beta-closed}
If \(G\in\mathrm{GL}_n\) and
\(\alpha=v^\top G^{-1}u\ne0\), the rank-one line has a corank-one crossing at
\(\lambda_{\rm sing}=-\alpha^{-1}\), and
\begin{equation}
 \beta=
 \frac{\alpha^2}
 {\|G^{-1}u\|_2\,\|G^{-\top}v\|_2}.
 \label{eq:beta-closed}
\end{equation}
\end{proposition}

\begin{proof}
The determinant lemma gives
\(\det(G+\lambda uv^\top)=\det(G)(1+\lambda\alpha)\).  The right and left
null spaces at its root are spanned by \(G^{-1}u\) and \(G^{-\top}v\).
After normalization, both numerators in
\(\beta=|u_0^\top u|\,|v_0^\top v|\) equal \(|\alpha|\), which proves
\eqref{eq:beta-closed}.  The factor
\(I-\alpha^{-1}G^{-1}uv^\top\) is the identity minus a rank-one projector,
so the corank is one.
\end{proof}

\subsection{The complete certificate slice and its uniqueness}

Fix a dual minimizer and write
\[
 A^\star=G+\lambda^\star\Theta=U_r\Sigma_rV_r^\top,
 \qquad
 \ell=m-r,\quad k=n-r,
 \qquad
 C=U_0^\top\Theta V_0,\qquad \beta=\|C\|_*.
\]
By \eqref{eq:solution-face}, every exact certificate is
\begin{equation}
 \Phi^\star=U_rV_r^\top+U_0KV_0^\top,
 \qquad
 K\in\mathcal K(C,a):=
 \{K:\|K\|_2\le1,\ \langle C,K\rangle=-a\}.
 \label{eq:all-directions}
\end{equation}

\begin{theorem}[Feasibility and uniqueness of the certificate slice]
\label{thm:certificate-uniqueness}
The slice \(\mathcal K(C,a)\) is nonempty if and only if \(|a|\le\beta\).
If \(\ell k=0\), then \(C=0\), the kernel-block space consists only of the
empty zero matrix, and the slice is feasible and a singleton exactly when
\(a=0\).  Assume henceforth that \(\ell,k\ge1\).  If \(C=0\), feasibility
forces \(a=0\), and the slice is the entire, non-singleton spectral unit ball.
If \(C\ne0\), then:
\begin{enumerate}
 \item if \(|a|<\beta\), the slice is a singleton if and only if
 \(\ell k=1\);
 \item if \(|a|=\beta\), the slice is a singleton if and only if
 \(\operatorname{rank}C=\min\{\ell,k\}\).
\end{enumerate}
In all remaining feasible cases with \(\ell,k\ge1\), it contains infinitely
many matrices.
Consequently, at a transverse square corank-one crossing the exact certificate
is unique in both the strict and boundary regimes.
\end{theorem}

\begin{proof}
Feasibility is spectral--nuclear duality.  The zero-dimensional case is
immediate.  If \(\ell,k\ge1\) and \(C=0\), the slice is the whole spectral
unit ball.  Now let
\(r_C=\operatorname{rank}C\), let
\(C=P_C\Sigma_C Q_C^\top\ne0\) be a compact SVD, and put
\(S_C=P_CQ_C^\top\).  When
\(|a|<\beta\), the feasible matrix
\(K_0=-(a/\beta)S_C\) has spectral norm strictly below one.  If \(\ell k>1\),
choose nonzero \(H\perp C\); both \(K_0\pm\varepsilon H\) remain feasible for
small \(\varepsilon\).  If \(\ell k=1\), the equality fixes the only entry.

At \(|a|=\beta\), equality in spectral--nuclear duality gives the complete
endpoint face
\begin{equation}
 \mathcal K(C,a)=
 \left\{-\operatorname{sign}(a)
 \bigl(P_CQ_C^\top+P_\perp WQ_\perp^\top\bigr):\|W\|_2\le1\right\},
 \label{eq:endpoint-face}
\end{equation}
where \(W\in\mathbb R^{(\ell-r_C)\times(k-r_C)}\), and
\(P_\perp,Q_\perp\) complete the singular bases.  Indeed, equality for each
positive singular pair forces
\(Kq_i=-\operatorname{sign}(a)p_i\) and
\(K^\top p_i=-\operatorname{sign}(a)q_i\), which kills both cross blocks.  The face
is a singleton exactly when one complementary space is zero, equivalently
when \(\operatorname{rank}C=\min\{\ell,k\}\).
\end{proof}

Thus compact-polar failure is not an artifact of certificate nonuniqueness:
in the square corank-one instances governing \cref{sec:rates}, both the
best approximant and the exact certificate are unique, but the compact choice
\(K=0\) is still wrong whenever \(a\ne0\).

\subsection{Hard spectral completion and the no-clipping dichotomy}

The minimum-Frobenius member of the complete slice can always be computed from
one SVD and one monotone scalar equation.  This also supplies step 4 of the
exact interval oracle.

\begin{theorem}[Minimum-Frobenius spectral completion]
\label{thm:frobenius-water-filling}
Let \(C\in\mathbb R^{\ell\times k}\) and suppose
\(|a|\le\beta:=\|C\|_*\).  If \(C=0\), then \(a=0\) and the unique
minimum-Frobenius completion is \(K_F=0\).  Otherwise, let
\(C=P_C\operatorname{diag}(\sigma_1,\ldots,\sigma_{r_C})Q_C^\top\) be a compact
SVD, where \(r_C=\operatorname{rank}C\) and
\(\sigma_1\ge\cdots\ge\sigma_{r_C}>0\).  Define
\[
 g_F(\eta)=\sum_{i=1}^{r_C}\sigma_i\min\{\eta\sigma_i,1\}.
\]
For \(|a|<\beta\), there is a unique
\(\eta\in[0,1/\sigma_{r_C})\) with \(g_F(\eta)=|a|\); at the endpoint take
\(\eta=1/\sigma_{r_C}\).  With \(\operatorname{sign}(0)=0\), the unique member of
\(\mathcal K(C,a)\) of minimum Frobenius norm is
\begin{equation}
 K_F=-\operatorname{sign}(a)P_C
 \operatorname{diag}\bigl(\min\{\eta\sigma_i,1\}\bigr)Q_C^\top.
 \label{eq:KF-main}
\end{equation}
At \(|a|=\beta\), the multiplier parameter need not be unique, but the
displayed matrix is unique.
\end{theorem}

\begin{proof}
The case \(C=0\) is immediate.  Otherwise, projection of \(\mu C\) onto the spectral unit ball clips its singular values
at one.  The continuous function \(g_F\) increases strictly until it reaches
\(\beta\), so the stated multiplier exists.  With
\(\mu=-\operatorname{sign}(a)\eta\), this projection is \(K_F\) and satisfies
the equality.  For any other feasible \(L\), projection optimality and
\(\langle C,L-K_F\rangle=0\) give
\(\langle K_F,L-K_F\rangle\ge0\).  Hence
\[
 \|L\|_F^2\ge\|K_F\|_F^2+\|L-K_F\|_F^2,
\]
which proves optimality and uniqueness.
\end{proof}

For the rank-one normal \eqref{eq:rank-one-normal}, put
\(x=U_0^\top u\), \(y=V_0^\top v\).  Then
\(C=xy^\top\) and \(\beta=\|x\|_2\|y\|_2\).  If \(\beta>0\), no clipping
occurs, and
\begin{equation}
 K_F=K_{\rm can}
 =-\frac{a}{\beta}
 \frac{x}{\|x\|_2}\frac{y^\top}{\|y\|_2}.
 \label{eq:Kcan-main}
\end{equation}
If \(\beta=0\), dual optimality forces \(a=0\), and
\(K_F=K_{\rm can}=0\).

\begin{proposition}[No-clipping/partial-isometry dichotomy]
\label{prop:partial-isometry}
For fixed \(C\ne0\), the unconstrained minimum-Frobenius solution
\[
 K_{\rm lin}=-\frac{a}{\|C\|_F^2}C
\]
is spectrally feasible if and only if
\(|a|\le\|C\|_F^2/\|C\|_2\).  Moreover
\begin{equation}
 \frac{\|C\|_F^2}{\|C\|_2}\le\|C\|_*,
 \label{eq:pi-inequality}
\end{equation}
with equality if and only if all nonzero singular values of \(C\) coincide.
Equivalently, \(K_{\rm lin}\) is feasible for every admissible scalar
right-hand side if and only if \(C\) is a scalar multiple of a partial
isometry.
\end{proposition}

\begin{proof}
The first claim follows from
\(\|K_{\rm lin}\|_2=|a|\|C\|_2/\|C\|_F^2\).  If \(\sigma_1=\|C\|_2\), then
\[
 \|C\|_F^2=\sum_i\sigma_i^2
 \le\sigma_1\sum_i\sigma_i=\|C\|_2\|C\|_*,
\]
with equality exactly when each positive \(\sigma_i\) equals \(\sigma_1\).
\end{proof}

For a general normal, \eqref{eq:KF-main} replaces the rank-one formula by one
SVD and one scalar solve.  The smooth analogue below has a different spectral
profile; \cref{cor:center-agreement-main} gives the exact condition under
which the two centers agree.

\section{Fenchel recovery and center selection}
\label{sec:smoothing}

For \(\delta>0\), define
\begin{equation}
 \psi_\delta(s)=\sqrt{s^2+\delta^2}-\delta,
 \qquad
 \Psi_\delta(A)=\sum_{i=1}^{q}\psi_\delta(\sigma_i(A)),
 \qquad q=\min\{m,n\}.
 \label{eq:Psi-main}
\end{equation}
This is a pseudo-Huber smoothing of the nuclear norm.  The spectral Fenchel
pair, its gradient formula, and its Frobenius Lipschitz continuity are known
\citep{mustafi2026move,lobos2025smooth,feoktistov2026softsign}; we use that
pair after imposing the tangency equality to identify the constrained
regularized limit inside \eqref{eq:all-directions}.

The scalar conjugate is
\begin{equation}
 \psi_\delta^*(z)=
 \begin{cases}
 \delta(1-\sqrt{1-z^2}),&|z|\le1,\\
 +\infty,&|z|>1.
 \end{cases}
 \label{eq:scalar-conjugate-main}
\end{equation}
Indeed, for \(|z|<1\), stationarity in the conjugate gives
\(s=\delta z/\sqrt{1-z^2}\), and the endpoints follow by continuity.
Define
\begin{equation}
 R(\Phi)=
 \sum_{i=1}^{q}\left(1-\sqrt{1-\sigma_i(\Phi)^2}\right)
 +\iota_{\{\|\Phi\|_2\le1\}}(\Phi).
 \label{eq:R-main}
\end{equation}

\begin{lemma}[Conjugacy, strong convexity, and smooth gradient]
\label{lem:conjugacy}
The regularizer \(R\) is \(1\)-strongly convex in the Frobenius norm, and
\begin{equation}
 \Psi_\delta=(\delta R)^*.
 \label{eq:matrix-conjugacy-main}
\end{equation}
The function \(\Psi_\delta\) is \(C^\infty\) and strictly convex on the full
matrix space.  If \(A=U\operatorname{diag}(\sigma_i)V^\top\) is any full SVD,
then
\begin{equation}
 Z_\delta(A):=\nabla\Psi_\delta(A)
 =U\operatorname{diag}\left(
 \frac{\sigma_i}{\sqrt{\sigma_i^2+\delta^2}}
 \right)V^\top.
 \label{eq:Zdelta}
\end{equation}
Equivalently,
\begin{equation}
 Z_\delta(A)=
 \begin{cases}
 A(A^\top A+\delta^2I)^{-1/2},&m\ge n,\\
 (AA^\top+\delta^2I)^{-1/2}A,&m<n.
 \end{cases}
 \label{eq:Zdelta-coordinate-free}
\end{equation}
\end{lemma}

\begin{proof}
On \((-1,1)\), the even scalar function
\(\rho(t)=1-\sqrt{1-t^2}\) satisfies
\(\rho''(t)=(1-t^2)^{-3/2}\ge1\).  Hence
\(\rho(t)-t^2/2\), extended by \(+\infty\) outside \([-1,1]\), is proper,
closed, and convex.  Spectral lifting of absolutely symmetric closed convex functions
\citep{lewis1995convex} shows that
\(R(\Phi)-\|\Phi\|_F^2/2\) is proper, closed, and convex.  Thus \(R\) is
lower semicontinuous, and this is the asserted strong convexity.
The same lifting applied to \eqref{eq:scalar-conjugate-main} proves
\eqref{eq:matrix-conjugacy-main}.

For strict matrix convexity, use the symmetric dilation
\[
 \mathcal H(A)=\begin{pmatrix}0&A\\A^\top&0\end{pmatrix},
 \qquad g_\delta(t)=\sqrt{t^2+\delta^2}.
\]
Its spectrum consists of \(\pm\sigma_i(A)\) and \(|m-n|\) zeros, so
\[
 \Psi_\delta(A)
 =\frac12\operatorname{tr}g_\delta(\mathcal H(A))
 -\frac{m+n}{2}\delta.
\]
For a symmetric matrix, the trace-spectral Hessian has divided differences
\((g_\delta'(x)-g_\delta'(y))/(x-y)\), with the diagonal value
\(g_\delta''(x)\) \citep{lewis2001twice}.  Since
\(g_\delta'(t)=t/\sqrt{t^2+\delta^2}\) is strictly increasing, all these
coefficients are positive.  A nonzero matrix perturbation has a nonzero
dilation, so the quadratic Hessian form is positive.  This proves strict
convexity; functional calculus gives smoothness and
\eqref{eq:Zdelta}--\eqref{eq:Zdelta-coordinate-free}.
\end{proof}

\begin{theorem}[Unique smooth oracle and Fenchel center path]
\label{thm:center-path}
For every \(\Theta\ne0\), the scalar function
\begin{equation}
 d_\delta(\lambda)=\Psi_\delta(G+\lambda\Theta)
 \label{eq:smooth-dual-main}
\end{equation}
is coercive and strictly convex.  It has a unique minimizer
\(\lambda_\delta\), characterized by
\begin{equation}
 h_\delta(\lambda)
 :=\langle\Theta,Z_\delta(G+\lambda\Theta)\rangle=0.
 \label{eq:smooth-root-main}
\end{equation}
The direction
\begin{equation}
 \Phi_\delta=Z_\delta(G+\lambda_\delta\Theta)
 \label{eq:Phi-delta-main}
\end{equation}
is exactly tangent and satisfies \(\|\Phi_\delta\|_2<1\).

Moreover, \(\Phi_\delta\) is the unique solution of
\begin{equation}
 \max_{\substack{\|\Phi\|_2\le1\\
                  \langle\Theta,\Phi\rangle=0}}
 \{\langle G,\Phi\rangle-\delta R(\Phi)\}.
 \label{eq:regularized-primal}
\end{equation}
If \(\mathcal P^\star\) is the exact solution face in
\eqref{eq:solution-face}, then
\begin{equation}
 \Phi_\delta\longrightarrow
 \Phi_R^\star:=\argmin_{\Phi\in\mathcal P^\star}R(\Phi)
 \qquad(\delta\downarrow0).
 \label{eq:center-limit-main}
\end{equation}
The center is unique.  If the exact dual minimizer is unique, then
\(\lambda_\delta\to\lambda^\star\).
\end{theorem}

\begin{proof}
Strict convexity follows by restricting \(\Psi_\delta\) to the injective
affine line \(G+\lambda\Theta\).  Since
\begin{equation}
 0\le\|A\|_*-\Psi_\delta(A)\le q\delta,
 \label{eq:smoothing-uniform}
\end{equation}
the lower bound
\(d_\delta(\lambda)\ge
 |\lambda|\|\Theta\|_*-\|G\|_*-q\delta\) proves coercivity.  Differentiation
gives \eqref{eq:smooth-root-main}; the derivative of a differentiable,
strictly convex, coercive scalar function has one zero.  Each multiplier in
\eqref{eq:Zdelta} lies in \([0,1)\), and the root gives tangency, proving
the first part.

Conjugacy \eqref{eq:matrix-conjugacy-main} gives
\[
 \Psi_\delta(G+\lambda\Theta)
 =\max_{\|\Phi\|_2\le1}
 \{\langle G+\lambda\Theta,\Phi\rangle-\delta R(\Phi)\}.
\]
The feasible set of \eqref{eq:regularized-primal} is compact and contains
the relative Slater point zero.  Lagrange duality is therefore exact, while
strong convexity of \(R\) makes the primal optimizer unique.  Fenchel equality
identifies it with \eqref{eq:Phi-delta-main}.

Let \(p^\star\) denote the unregularized optimum.  Optimality against the
unique \(R\)-minimizer \(\Phi_R^\star\) on \(\mathcal P^\star\) gives
\begin{equation}
 0\le p^\star-\langle G,\Phi_\delta\rangle
 \le\delta\{R(\Phi_R^\star)-R(\Phi_\delta)\}.
 \label{eq:center-gap-main}
\end{equation}
In particular, \(R(\Phi_\delta)\le R(\Phi_R^\star)\) and, because
\(R\ge0\), the linear gap is at most
\(\delta R(\Phi_R^\star)\to0\).
Compactness makes every sequence \(\delta_k\downarrow0\) have a convergent
subsequence.  The linear gap in \eqref{eq:center-gap-main} puts every limit in
\(\mathcal P^\star\); lower semicontinuity gives an \(R\)-value no larger than
that of \(\Phi_R^\star\).  Strong convexity makes the center unique, so all
subsequences have the same limit.

 Finally, the uniform convergence
 \eqref{eq:smoothing-uniform} and common coercive lower bound imply convergence
 of the scalar minimizers whenever the exact minimizer is unique.
\end{proof}

The abstract center has an explicit spectral form for every matrix normal.
It is generally different from the hard minimum-Frobenius completion
\eqref{eq:KF-main}.

\begin{theorem}[General pseudo-Huber certificate center]
\label{thm:general-fenchel-center}
At a rank-deficient optimum, use \(C\), \(a\), and \(\beta\) from
\eqref{eq:a-C-beta-main}.  If \(C=0\), then \(a=0\) and \(K_R=0\).
Otherwise, let
\(C=P_C\operatorname{diag}(\sigma_1,\ldots,\sigma_{r_C})Q_C^\top\) be a compact
SVD, where \(r_C=\operatorname{rank}C\).  If \(|a|<\beta\), there is a unique
\(\eta\ge0\) satisfying
\begin{equation}
 \sum_{i=1}^{r_C}
 \frac{\eta\sigma_i^2}{\sqrt{1+\eta^2\sigma_i^2}}=|a|.
 \label{eq:eta-R-main}
\end{equation}
The kernel block of \(\Phi_R^\star\) is
\begin{equation}
 K_R=-\operatorname{sign}(a)P_C
 \operatorname{diag}\left(
  \frac{\eta\sigma_i}{\sqrt{1+\eta^2\sigma_i^2}}
 \right)Q_C^\top,
 \label{eq:KR-main}
\end{equation}
with \(K_R=0\) when \(a=0\).  At \(|a|=\beta\), the formula is understood as
\(\eta\to\infty\), giving
\(K_R=-\operatorname{sign}(a)P_CQ_C^\top\).  For rank-one \(C\),
\eqref{eq:KR-main} reduces exactly to
\(K_{\rm can}\) in \eqref{eq:Kcan-main}.
\end{theorem}

\begin{proof}
In the singular bases,
\(R(\Phi)=\operatorname{rank}(A^\star)+\mathcal R_0(K)\), where
\[
 \mathcal R_0(K)=\sum_j\rho(\sigma_j(K))
 +\iota_{\{\|K\|_2\le1\}}(K),
 \qquad \rho(t)=1-\sqrt{1-t^2}.
\]
The scalar conjugate is \(\rho^*(z)=\sqrt{1+z^2}-1\).  Spectral lifting gives
\[
 \mathcal R_0^*(Y)=
 \sum_{j=1}^{\min\{\ell,k\}}
 \bigl(\sqrt{1+\sigma_j(Y)^2}-1\bigr).
\]
For \(|a|<\beta\), the feasible point
\(-(a/\beta)P_CQ_C^\top\) has spectral norm below one, so Slater holds.  The
KKT condition is
\(K_R=\nabla\mathcal R_0^*(\mu C)\), with
\(\mu=-\operatorname{sign}(a)\eta\).  It gives \eqref{eq:KR-main} and
\[
 \langle C,K_R\rangle
 =-\operatorname{sign}(a)\sum_{i=1}^{r_C}
 \frac{\eta\sigma_i^2}{\sqrt{1+\eta^2\sigma_i^2}}=-a.
\]
The left side
of \eqref{eq:eta-R-main} starts at zero, has derivative
\[
 \sum_i\frac{\sigma_i^2}{(1+\eta^2\sigma_i^2)^{3/2}}>0,
\]
and tends to \(\sum_i\sigma_i=\beta\), proving existence and uniqueness for
the strict slice.  At the endpoint, \eqref{eq:endpoint-face} fixes the support
block and leaves only \(W\); nonnegativity of \(\rho\), with equality only at
zero, selects \(W=0\).  The case \(C=0\) is the same unconstrained minimum.
For rank-one \(C\), the spectral center formula has one active coefficient,
and \eqref{eq:eta-R-main} fixes it at \(|a|/\beta\), yielding
\eqref{eq:Kcan-main}.
\end{proof}

\begin{corollary}[Exact agreement of the hard and smooth centers]
\label{cor:center-agreement-main}
Let \(C\ne0\) and \(|a|\le\beta\).  Then \(K_F=K_R\) if and only if at least
one of the following holds: \(a=0\), \(|a|=\beta\), or all positive singular
values of \(C\) coincide.  Thus, in the strict nonzero regime, the two centers
agree exactly for scalar multiples of partial isometries.
\end{corollary}

\begin{proof}
Agreement is immediate at zero and at the endpoint.  If all positive singular
values are equal, symmetry makes every active coefficient in both formulas
equal to \(|a|/\beta\).  Conversely, suppose
\(0<|a|<\beta\) and the centers agree.  Every coefficient in the smooth
formula is strictly below one, so the hard formula is unsaturated.  Equality
of the coefficients gives, for every positive \(\sigma_i\),
\[
 \eta_F=\frac{\eta_R}{\sqrt{1+\eta_R^2\sigma_i^2}}.
\]
Since \(\eta_R>0\), all \(\sigma_i\) must coincide.
\end{proof}

\subsection{A smoothing family and two distinguished centers}

The preceding center is not an isolated artifact of pseudo-Huber smoothing.
Call a scalar profile \(\psi\) \emph{admissible} if it is even, convex, and
\(C^1\), with \(\psi(0)=0\), if \(0\le\psi'(x)\le1\) for \(x\ge0\), and if
\(\psi'\) is nondecreasing with \(\lim_{x\to\infty}\psi'(x)=1\).  We further require
\[
 c_\psi:=\lim_{x\to\infty}\{x-\psi(x)\}<\infty,
 \qquad
 0\le x<y,\ \psi'(y)<1\ \Longrightarrow\ \psi'(x)<\psi'(y),
 \tag{A}
 \label{eq:admissible-profile-main}
\]
and assume that, for some \(\mu_\psi>0\), the extended-value function
\(z\mapsto\psi^*(z)-\mu_\psi z^2/2\) is convex on \(\mathbb R\); the domain
of \(\psi^*\) is \([-1,1]\).  The second condition in
\eqref{eq:admissible-profile-main} means strict increase before possible
finite saturation; it includes both profiles below.  Define
\begin{equation}
 \Psi_\delta^\psi(A)=\delta\sum_{i=1}^q
 \psi\!\left(\frac{\sigma_i(A)}{\delta}\right),
 \qquad
 R_\psi(\Phi)=\sum_{i=1}^q\psi^*(\sigma_i(\Phi)).
 \label{eq:profile-lifts-main}
\end{equation}
The last expression is understood as \(+\infty\) when
\(\|\Phi\|_2>1\).

\begin{proposition}[Admissible profiles select spectral centers]
\label{prop:smoothing-family-main}
For every admissible \(\psi\),
\(\Psi_\delta^\psi=(\delta R_\psi)^*\), and \(R_\psi\) is strongly convex
in the Frobenius norm.  If \(\lambda_\delta^\psi\) is any minimizer of
\(\Psi_\delta^\psi(G+\lambda\Theta)\), then
\[
 \Phi_\delta^\psi=\nabla\Psi_\delta^\psi
   (G+\lambda_\delta^\psi\Theta)
\]
is independent of the choice of multiplier and is the unique maximizer of
\eqref{eq:regularized-primal} with \(R\) replaced by \(R_\psi\).  Moreover,
\begin{equation}
 \Phi_\delta^\psi\longrightarrow
 \Phi_\psi^\star:=\argmin_{\Phi\in\mathcal P^\star}R_\psi(\Phi).
 \label{eq:profile-center-limit-main}
\end{equation}

At a strict rank-deficient minimizer, with \(a,C,\beta\) as in
\eqref{eq:a-C-beta-main}, every choice of \(\lambda_\delta^\psi\) satisfies
\begin{equation}
 \frac{\lambda_\delta^\psi-\lambda^\star}{\delta}
 \longrightarrow t_0^\psi=-\sign(a)\eta_\psi,
 \qquad
 \sum_{i=1}^{r_C}\sigma_i\psi'(\eta_\psi\sigma_i)=|a|,
 \label{eq:profile-strict-rate-main}
\end{equation}
where \(\eta_\psi\ge0\) is the unique solution.  The null--null block of
the center is
\begin{equation}
 K_\psi=-\sign(a)P_C\diag\!\bigl(
 \psi'(\eta_\psi\sigma_i)\bigr)Q_C^\top.
 \label{eq:profile-center-block-main}
\end{equation}
The same sharp-growth argument gives the uniform localization
\begin{equation}
 |\lambda_\delta^\psi-\lambda^\star|
 \le \frac{q c_\psi}{\beta-|a|}\,\delta.
 \label{eq:profile-strict-bound-main}
\end{equation}
For
\begin{equation}
 \psi_{\rm ph}(x)=\sqrt{1+x^2}-1,
 \qquad
 \psi_{\rm hub}(x)=
 \begin{cases}x^2/2,&|x|\le1,\\ |x|-1/2,&|x|>1,
 \end{cases}
 \label{eq:two-profile-members-main}
\end{equation}
these formulas give, respectively, \((\eta_R,K_R)\) and
\((\eta_F,K_F)\).  In particular,
\begin{equation}
 R_{\rm hub}(\Phi)=\tfrac12\|\Phi\|_F^2
 +\iota_{\{\|\Phi\|_2\le1\}}(\Phi),
 \label{eq:huber-regularizer-main}
\end{equation}
so the hard water-filled completion is itself a Fenchel center.
\end{proposition}

\begin{proof}
Scalar perspective conjugacy and spectral lifting give
\(\Psi_\delta^\psi=(\delta R_\psi)^*\); strong convexity also lifts from
the absolutely symmetric singular-value function.  Hence the proof of
\cref{thm:center-path}, without its strict-convexity claim for the scalar
multiplier, gives the unique regularized certificate and
\eqref{eq:profile-center-limit-main}.  The uniform bound
\(0\le |x|-\delta\psi(x/\delta)\le c_\psi\delta\) replaces
\eqref{eq:smoothing-uniform}.

For the local rate, use the exact analytic cluster frames from
\cref{thm:direction-strict-general}.  After setting
\(s=\delta t\), the separated positive cluster contributes \(at+o(1)\),
whereas the small block contributes
\(\sum_i\psi(t\sigma_i)+o(1)\), uniformly for bounded \(t\).  Thus the
rescaled objectives converge compact-uniformly to
\[
 F_\psi(t)=at+\sum_{i=1}^{r_C}\psi(t\sigma_i).
\]
It is coercive because \(|a|<\beta\).  The map
\(g_\psi(t)=F_\psi'(t)-a\) is strictly increasing wherever
\(|g_\psi(t)|<\beta\) by \eqref{eq:admissible-profile-main}.  Since
\(|a|<\beta\), \(F_\psi'\) has the unique zero in
\eqref{eq:profile-strict-rate-main}.  The sharp-growth localization used in
\cref{thm:rate-strict-general} then gives convergence of every minimizing
multiplier.  Finally, KKT conditions for minimizing the null-block part of
\(R_\psi\) over \(\mathcal K(C,a)\), or equivalently differentiating its
spectral conjugate at \(-\sign(a)\eta_\psi C\), give
\eqref{eq:profile-center-block-main}.

For pseudo-Huber, \(\psi'_{\rm ph}(x)=x/\sqrt{1+x^2}\), recovering
\eqref{eq:eta-R-main}--\eqref{eq:KR-main}.  For Huber,
\(\psi'_{\rm hub}(x)=\min\{x,1\}\) on \([0,\infty)\) and
\(\psi_{\rm hub}^*(z)=z^2/2\) on \([-1,1]\), which gives
\eqref{eq:huber-regularizer-main}; its scalar equation is exactly the hard
water-filling equation defining \(\eta_F\) and \(K_F\).
\end{proof}

\begin{remark}[Variational center versus analytic oracle]
\label{rem:huber-caveat-main}
Huber smoothing is \(C^1\), convex, and coercive, so its scalar stationarity
equation always has a solution.  Its derivative saturates, however, so the
objective need not be strictly convex and can have locally constant
multiplier intervals; the clipping map is also nonanalytic at its threshold.
Thus multiplier uniqueness and the analytic direction expansion of
\cref{thm:direction-strict-general} do not transfer without extra active-set
nondegeneracy.  Huber is used here to identify \(K_F\) variationally and to
separate center selection from the tail mechanism governing boundary rates,
not to replace the pseudo-Huber oracle.
For example, if \(G=2I_2\), \(\Theta=\diag(1,-1)\), and \(0<\delta<2\),
then the full Huber multiplier interval is
\([-2+\delta,2-\delta]\), while every multiplier produces the same
certificate \(I_2\).
\end{remark}

\begin{remark}[Approximant versus certificate regularization]
\label{rem:legg-ward}
The canonical trace-class approximation of \citet{legg1985canonical} uses the
limit of Schatten \(c_p\)-approximations as \(p\downarrow1\) to resolve
possible nonuniqueness on the approximant side.  Our regularization acts on
the certificate side under the additional tangency equality.  In the no-root
regime the approximant is unique by \cref{thm:interval}, whereas the
certificate may or may not be unique by
\cref{thm:certificate-uniqueness}.  Thus \cref{thm:center-path} is recovery in
the singleton case and genuine Fenchel-center selection otherwise; the rate
theory below applies to the corresponding regularized certificate path.
\end{remark}

\begin{corollary}[Global safety bounds]
\label{cor:global-bounds}
For every \(A\), \eqref{eq:smoothing-uniform} holds, and
\begin{equation}
 0\le p^\star-\langle G,\Phi_\delta\rangle\le q\delta,
 \qquad
 |\lambda_\delta|
 \le\frac{2\|G\|_*+q\delta}{\|\Theta\|_*}.
 \label{eq:smooth-lambda-bound}
\end{equation}
\end{corollary}

\begin{proof}
For \(s\ge0\),
\(0\le s-(\sqrt{s^2+\delta^2}-\delta)\le\delta\); summing proves
\eqref{eq:smoothing-uniform}.  Since \(R\le q\) on the spectral unit ball,
\eqref{eq:center-gap-main} gives the primal bias.  Finally,
\(d_\delta(\lambda_\delta)\le d_\delta(0)\le\|G\|_*\), whereas the coercive
lower bound used above gives the multiplier estimate.
\end{proof}

The exact interval method needs no smoothing if one merely wants an exact
direction.  The value of \eqref{eq:smooth-root-main} is different: it is a
continuous strictly monotone equation, avoids an exact-rank branch, and
recovers or selects the explicit variational center.  The next section quantifies its
behavior at rank loss.

\section{Local smoothing laws at regular and singular minimizers}
\label{sec:rates}

Write
\[
 A(\lambda)=G+\lambda\Theta,
 \qquad s=\lambda-\lambda^\star.
\]
All statements in this section are local asymptotics at a fixed instance
under the printed nondegeneracy assumptions.  The strict rank-deficient law
below permits arbitrary shape and corank.  Only the boundary and crossover
laws use a transverse square corank-one crossing, because their proof isolates
one signed Rellich-analytic small mode.

\subsection{Full-rank regular minimizers}

\begin{theorem}[Quadratic regular displacement]
\label{thm:rate-regular}
Assume that \(\lambda^\star\) minimizes \(\phi\), that
\(A^\star=A(\lambda^\star)\) has full rectangular rank, and that
\(\kappa=\phi''(\lambda^\star)>0\).  If
\(A^\star=U\Sigma V^\top\), then
\begin{equation}
 \lambda_\delta-\lambda^\star
 =-\frac{b^\star}{\kappa}\delta^2+o(\delta^2),
 \qquad
 b^\star=-\frac12
 \langle\Theta,U\Sigma^{-2}V^\top\rangle.
 \label{eq:regular-rate-main}
\end{equation}
The matrix in the pairing is basis independent:
\begin{equation}
 U\Sigma^{-2}V^\top=
 \begin{cases}
 A^\star((A^\star)^\top A^\star)^{-3/2},&m\ge n,\\
 (A^\star(A^\star)^\top)^{-3/2}A^\star,&m<n.
 \end{cases}
 \label{eq:regular-intrinsic-main}
\end{equation}
Thus the displacement is \(O(\delta^2)\), and it is
\(\Theta(\delta^2)\) when \(b^\star\ne0\).
\end{theorem}

\begin{proof}
In a full-rank neighborhood,
\[
 \frac{\sigma}{\sqrt{\sigma^2+\delta^2}}
 =1-\frac{\delta^2}{2\sigma^2}+O(\delta^4)
\]
uniformly.  Functional calculus therefore gives
\[
 Z_\delta(A(\lambda))
 =\Sc(A(\lambda))
 -\frac{\delta^2}{2}
 U(\lambda)\Sigma(\lambda)^{-2}V(\lambda)^\top
 +O(\delta^4).
\]
The displayed matrix functions are analytic and basis independent even at
repeated positive singular values.  Pairing with \(\Theta\) yields
\(h_\delta(\lambda)=\phi'(\lambda)+\delta^2b(\lambda)+O(\delta^4)\), with
\(b(\lambda^\star)=b^\star\).  The scalar implicit function theorem at
\(\phi'(\lambda^\star)=0\), \(\phi''(\lambda^\star)=\kappa\) proves
\eqref{eq:regular-rate-main}.
\end{proof}

\subsection{Strict rank-deficient displacement at arbitrary corank}

\begin{theorem}[Strict displacement at arbitrary corank]
\label{thm:rate-strict-general}
Let \(\lambda^\star\) minimize \(\phi\), let
\(A^\star=A(\lambda^\star)\) be rank deficient, and define \(a,C,\beta\)
by \eqref{eq:a-C-beta-main}.  Assume the strict condition
\(|a|<\beta\).  Let
\(\sigma_1\ge\cdots\ge\sigma_{r_C}>0\) be the positive singular values
of \(C\), and let \(\eta_R\ge0\) be the unique solution of
\begin{equation}
 \sum_{i=1}^{r_C}
 \frac{\eta_R\sigma_i^2}{\sqrt{1+\eta_R^2\sigma_i^2}}=|a|.
 \label{eq:eta-strict-rate-main}
\end{equation}
Then
\begin{equation}
 \frac{\lambda_\delta-\lambda^\star}{\delta}
 \longrightarrow-\sign(a)\eta_R.
 \label{eq:strict-general-rate-main}
\end{equation}
Moreover, with \(q=\min\{m,n\}\),
\begin{equation}
 |\lambda_\delta-\lambda^\star|
 \le \frac{q}{\beta-|a|}\,\delta.
 \label{eq:strict-general-bound-main}
\end{equation}
No square-shape, corank-one, or rank-one-normal assumption is required.
The displacement is \(\Theta(\delta)\) exactly when \(a\ne0\); when
\(a=0\), it is \(o(\delta)\).
\end{theorem}

\begin{proof}
Put \(q=\min\{m,n\}\), \(r=\rank A^\star\), and
\(t_\delta=(\lambda_\delta-\lambda^\star)/\delta\).  Strict optimality
and \cref{thm:interval} give the global sharp-growth estimate
\begin{equation}
 \phi(\lambda)\ge \phi(\lambda^\star)
 +(\beta-|a|)|\lambda-\lambda^\star|.
 \label{eq:strict-sharp-growth-main}
\end{equation}
Combining it with \eqref{eq:smoothing-uniform} and optimality of
\(\lambda_\delta\) yields \eqref{eq:strict-general-bound-main}.
Thus \((t_\delta)\) is bounded before any local expansion is used.

Define the rescaled convex functions
\[
 F_\delta(t)=
 \frac{\Psi_\delta(A(\lambda^\star+\delta t))
       -\Psi_\delta(A^\star)}{\delta}.
\]
Choose a spectral gap separating the \(r\) positive singular values of
\(A^\star\) from zero.  The sum \(\ell(s)\) of their analytic separated
cluster is smooth even when singular values inside that cluster cross, and
\(\ell'(0)=\langle\Theta,U_rV_r^\top\rangle=a\).  Since
\(\psi_\delta(\sigma)=\sigma-\delta+O(\delta^2)\) uniformly on the
separated cluster, its contribution to \(F_\delta(t)\) converges to \(at\),
uniformly for \(t\) in compact sets.

For the small cluster, let \(\Pi_V(s),\Pi_U(s)\) be the Riesz projections of
\(A(s)^\top A(s)\) and \(A(s)A(s)^\top\) onto the clusters converging to
zero, where \(A(s)=A(\lambda^\star+s)\).  The fixed gap makes them analytic,
with analytic orthonormal frames \(V_0(s),U_0(s)\)
\citep[Chapter~II]{kato1995perturbation}.  Resolvent intertwining gives
\(A(s)\Pi_V(s)=\Pi_U(s)A(s)\) and its transpose identity, so the cluster
subspaces reduce \(A(s)\) exactly.  Its small singular values are therefore
precisely those of
\[
 B(s)=U_0(s)^\top A(\lambda^\star+s)V_0(s)=sC+O(s^2)
\]
because \(B(0)=0\) and \(B'(0)=U_0^\top\Theta V_0=C\).  Weyl's
singular-value bound now gives, uniformly near zero,
\[
 \sigma_i(A(\lambda^\star+s))
 =|s|\sigma_i(C)+O(s^2),
\]
where zero singular values of \(C\) are included.  Setting \(s=\delta t\)
and using that \(x\mapsto\sqrt{1+x^2}-1\) is Lipschitz gives compact-uniform
convergence
\begin{equation}
 F_\delta(t)\longrightarrow
 F(t):=at+\sum_{i=1}^{r_C}
 \left(\sqrt{1+\sigma_i^2t^2}-1\right).
 \label{eq:strict-epi-limit-main}
\end{equation}

The function \(F\) is strictly convex, and
\(F(t)=at+\beta|t|+o(|t|)\); hence \(|a|<\beta\) makes it coercive.
Its unique minimizer solves
\[
 a+\sum_{i=1}^{r_C}
 \frac{\sigma_i^2t}{\sqrt{1+\sigma_i^2t^2}}=0,
\]
so it equals \(-\sign(a)\eta_R\).  The localization
\eqref{eq:strict-general-bound-main} puts every minimizer of \(F_\delta\) in
one fixed compact interval.  Compact-uniform convergence then implies that
every accumulation point minimizes \(F\); uniqueness proves
\eqref{eq:strict-general-rate-main}.  The final order statements follow
because \(\eta_R>0\) if and only if \(a\ne0\).
\end{proof}

\begin{remark}[One scalar governs both path and limit]
\label{rem:one-scalar-path-center}
Throughout this section the subscript \(R\) marks objects attached to the
pseudo-Huber regularizer, rather than a common tensor type.
The multiplier \(\eta_R\) in \eqref{eq:strict-general-rate-main} is exactly
the multiplier in the smooth certificate formula
\eqref{eq:eta-R-main}--\eqref{eq:KR-main}.  Thus the first displacement of
the smoothing path and the kernel block of its limiting certificate are two
manifestations of the same scalar balance.  More precisely,
\(K_R=Z_1(t_0C)\), where
\(t_0=-\operatorname{sign}(a)\eta_R\); the next theorem shows that this
identity is the zero-scale value of an analytic blow-up path.
\end{remark}

\begin{theorem}[First-order direction law at a strict minimizer]
\label{thm:direction-strict-general}
Under the hypotheses of \cref{thm:rate-strict-general}, put
\(t_0=-\operatorname{sign}(a)\eta_R\).  For the fixed instance there is a
jointly real-analytic blow-up extension \(\mathcal Z(\varepsilon,t)\), defined
near \((0,t_0)\), such that, for \(\varepsilon>0\),
\begin{equation}
 \mathcal Z(\varepsilon,t)
 =Z_\varepsilon(A(\lambda^\star+\varepsilon t)),
 \qquad
 \mathcal Z(0,t)=U_rV_r^\top+U_0Z_1(tC)V_0^\top.
 \label{eq:strict-blowup-map-main}
\end{equation}
Let
\[
 H(\varepsilon,t)=\langle\Theta,\mathcal Z(\varepsilon,t)\rangle,
 \qquad
 d_R:=\partial_tH(0,t_0)
 =\sum_{i=1}^{r_C}
 \frac{\sigma_i^2}{(1+t_0^2\sigma_i^2)^{3/2}}>0,
\]
and set
\begin{equation}
 t_1=-\frac{\partial_\varepsilon H(0,t_0)}{d_R},
 \qquad
 \mathcal M_R=\partial_\varepsilon\mathcal Z(0,t_0)
 -\frac{\langle\Theta,\partial_\varepsilon\mathcal Z(0,t_0)\rangle}{d_R}
  \partial_t\mathcal Z(0,t_0).
 \label{eq:strict-direction-constant-main}
\end{equation}
Then
\begin{align}
 \lambda_\delta-\lambda^\star
 &=\delta t_0+\delta^2t_1+O(\delta^3),
 \label{eq:strict-multiplier-second-main}\\
 \Phi_\delta
 &=\Phi_R^\star+\delta\mathcal M_R+O(\delta^2).
 \label{eq:strict-direction-first-main}
\end{align}
Consequently the returned-direction error is \(O(\delta)\), and it is
asymptotic to \(\delta\|\mathcal M_R\|_F\) when
\(\mathcal M_R\ne0\).  If \(\mathcal M_R=0\), it is \(O(\delta^2)\).
In particular, \(a=0\) forces \(t_0=t_1=0\) and
\(\mathcal M_R=0\), so that
\begin{equation}
 \lambda_\delta-\lambda^\star=O(\delta^3),
 \qquad
 \|\Phi_\delta-\Phi_R^\star\|_F=O(\delta^2).
 \label{eq:strict-zero-superconvergence-main}
\end{equation}
The constants are local to the fixed instance; no uniformity as the spectral
gap or \(\beta-|a|\) tends to zero is asserted.
\end{theorem}

\begin{proof}
Let \(r=\operatorname{rank}A^\star\).  The gap between its positive and zero
singular-value clusters gives local analytic orthogonal frames
\(U(s)=[U_+(s)\ U_0(s)]\) and \(V(s)=[V_+(s)\ V_0(s)]\) for which
\[
 U(s)^\top A(\lambda^\star+s)V(s)
 =\begin{pmatrix}A_+(s)&0\\0&B(s)\end{pmatrix}.
\]
Here \(A_+(s)\in\mathbb R^{r\times r}\) is invertible,
\(B(s)\in\mathbb R^{(m-r)\times(n-r)}\), and
\(B(0)=0\), \(B'(0)=C\).  Analytic division gives
\(B(s)=s\widetilde B(s)\), with \(\widetilde B(0)=C\); unequal left and
right nullities and zero singular values of \(C\) cause no difficulty.
The frames are chosen to agree at \(s=0\) with the fixed compact and null
frames used to define \(a\) and \(C\).

Write
\[
 \mathcal E_s(K)=U_0(s)KV_0(s)^\top,
 \qquad
 P_\varepsilon(s)=
 U_+(s)Z_\varepsilon(A_+(s))V_+(s)^\top.
\]
Orthogonal equivariance, block separation, and the scaling identity
\(Z_\varepsilon(B)=Z_1(B/\varepsilon)\) give, for \(\varepsilon>0\),
\begin{equation}
 Z_\varepsilon(A(\lambda^\star+\varepsilon t))
 =P_\varepsilon(\varepsilon t)
 +\mathcal E_{\varepsilon t}
  \left[Z_1\!\left(t\widetilde B(\varepsilon t)\right)\right].
 \label{eq:strict-blowup-block-main}
\end{equation}
The positive block is analytic in \((s,\varepsilon^2)\), while \(Z_1\) is
real analytic on the entire rectangular matrix space.  Thus the right-hand
side extends jointly analytically across \(\varepsilon=0\) and proves
\eqref{eq:strict-blowup-map-main}.  This construction uses analytic cluster
subspaces, not analytic individual singular vectors, so repeated positive
singular values are allowed \citep[Chapter II]{kato1995perturbation}.

At zero scale,
\[
 H(0,t)=a+\langle C,Z_1(tC)\rangle
 =a+\sum_{i=1}^{r_C}
   \frac{t\sigma_i^2}{\sqrt{1+t^2\sigma_i^2}}.
\]
Its unique zero is \(t_0\), and differentiation gives the stated positive
constant \(d_R\).  The analytic implicit function theorem therefore gives
\(t(\varepsilon)=t_0+\varepsilon t_1+O(\varepsilon^2)\).  For positive
\(\varepsilon\), strict convexity and
\cref{thm:rate-strict-general} identify this local branch with
\((\lambda_\varepsilon-\lambda^\star)/\varepsilon\).
Since \(Z_1(t_0C)=K_R\), Taylor expansion of
\(\mathcal Z(\delta,t(\delta))\) proves
\eqref{eq:strict-multiplier-second-main}--\eqref{eq:strict-direction-first-main}
and \eqref{eq:strict-direction-constant-main}.  When \(a=0\), one has
\(t_0=0\), while
\(Z_\varepsilon(A^\star)=\Sc(A^\star)+O(\varepsilon^2)\).  Hence
\(\partial_\varepsilon H(0,0)=0\) and
\(\partial_\varepsilon\mathcal Z(0,0)=0\), proving
\eqref{eq:strict-zero-superconvergence-main}.
\end{proof}

\begin{proposition}[Intrinsic computable first-correction formula]
\label{prop:strict-constants-closed-main}
Use the analytic cluster frames in the proof of
\cref{thm:direction-strict-general}, write
\(\widetilde B(s)=B(s)/s\) with its analytic value
\(\widetilde B(0)=C\), and define the basis-free frozen-scale field
\begin{equation}
 \mathcal W(s,t)=P_{\rm big}(s)
 +U_0(s)Z_1\!\bigl(t\widetilde B(s)\bigr)V_0(s)^\top,
 \qquad
 P_{\rm big}(s)=U_+(s)\Sc(A_+(s))V_+(s)^\top.
 \label{eq:frozen-field-main}
\end{equation}
Then \(\mathcal W(0,t_0)=\Phi_R^\star\) and
\begin{equation}
 \partial_\varepsilon\mathcal Z(0,t_0)
 =t_0\partial_s\mathcal W(0,t_0).
 \label{eq:frozen-chain-main}
\end{equation}
Consequently the constants in
\eqref{eq:strict-direction-constant-main} are
\begin{align}
 t_1&=-\frac{t_0
 \langle\Theta,\partial_s\mathcal W(0,t_0)\rangle}{d_R},
 \label{eq:t1-closed-main}\\
 \mathcal M_R&=t_0\left[
 \partial_s\mathcal W(0,t_0)
 -\frac{\langle\Theta,\partial_s\mathcal W(0,t_0)\rangle}{d_R}
  \partial_t\mathcal W(0,t_0)
 \right].
 \label{eq:MR-closed-main}
\end{align}
In particular, \(t_0=0\) makes both constants vanish.  If \(t_0\ne0\),
then
\begin{align*}
 \mathcal M_R=0
 &\iff \partial_s\mathcal W(0,t_0)
       \in\operatorname{span}\{\partial_t\mathcal W(0,t_0)\},\\
 t_1=0
 &\iff \langle\Theta,\partial_s\mathcal W(0,t_0)\rangle=0,
\end{align*}
and both vanish if and only if
\(\partial_s\mathcal W(0,t_0)=0\).
\end{proposition}

\begin{proof}
Both terms in \eqref{eq:frozen-field-main} are invariant under analytic
orthogonal changes of the left and right cluster frames, by orthogonal
equivariance of \(Z_1\).  In \eqref{eq:strict-blowup-block-main}, the positive
block satisfies
\(P_\varepsilon(s)=P_{\rm big}(s)+O(\varepsilon^2)\), jointly analytically
in \((s,\varepsilon^2)\).  All remaining \(\varepsilon\)-dependence enters
through \(s=\varepsilon t\).  Hence
\[
 \mathcal Z(\varepsilon,t)
 =\mathcal W(\varepsilon t,t)+O(\varepsilon^2),
\]
and the chain rule proves \eqref{eq:frozen-chain-main}.  Substitution into
\eqref{eq:strict-direction-constant-main} gives
\eqref{eq:t1-closed-main}--\eqref{eq:MR-closed-main}.  Finally,
\(\langle\Theta,\partial_t\mathcal W(0,t_0)\rangle=d_R>0\);
the stated vanishing alternatives follow directly.  Notice in particular
that parallelism alone cancels \(\mathcal M_R\), but need not cancel \(t_1\).
The full-rank polar-block derivatives in this formula may be evaluated by
standard Fr\'echet-derivative algorithms \citep{gawlik2017frechet}; the new
feature here is their coupling to the analytic rank-deficient blow-up.
\end{proof}

\begin{corollary}[Sharp compact-root superconvergence]
\label{cor:compact-root-superconvergence-main}
Under \cref{thm:direction-strict-general}, suppose \(a=0\).  By
\cref{thm:interval}, this is exactly compact-polar stationarity at
\(\lambda^\star\).  Let
\begin{equation}
 J_\star=U_r\Sigma_r^{-2}V_r^\top,
 \qquad d_0=\|C\|_F^2,
 \qquad \chi_\star=\langle\Theta,J_\star\rangle.
 \label{eq:compact-root-coefficients-main}
\end{equation}
where \(J_\star=0\) when the positive singular block is empty.  Strictness
gives \(C\ne0\), hence \(d_0>0\).
Then \(\Phi_R^\star=\Sc(A^\star)\), and
\begin{align}
 \lambda_\delta-\lambda^\star
 &=\frac{\chi_\star}{2d_0}\,\delta^3+O(\delta^4),
 \label{eq:compact-root-cubic-main}\\
 \Phi_\delta-\Phi_R^\star
 &=\delta^2\mathcal N_0+O(\delta^3),
 \qquad
 \mathcal N_0=-\frac12J_\star
 +\frac{\chi_\star}{2d_0}U_0CV_0^\top.
 \label{eq:compact-root-quadratic-main}
\end{align}
Moreover,
\begin{equation}
 \|\mathcal N_0\|_F^2
 =\frac14\|J_\star\|_F^2+\frac{\chi_\star^2}{4d_0}.
 \label{eq:compact-root-direction-norm-main}
\end{equation}
Thus the multiplier error is exactly cubic when \(\chi_\star\ne0\), and the
direction error is exactly quadratic whenever \(\rank A^\star>0\).  If
\(\chi_\star=0\), put
\begin{equation}
 J_2=U_r\Sigma_r^{-4}V_r^\top,
 \qquad \chi_2=\langle\Theta,J_2\rangle .
 \label{eq:compact-root-second-coefficients-main}
\end{equation}
Then the next coefficients are
\begin{align}
 \lambda_\delta-\lambda^\star
 &=-\frac{3\chi_2}{8d_0}\,\delta^5+O(\delta^6),
 \label{eq:compact-root-quintic-main}\\
 \Phi_\delta-\Phi_R^\star
 &=-\frac{\delta^2}{2}J_\star
 +\delta^4\left[\frac38J_2
 -\frac{3\chi_2}{8d_0}U_0CV_0^\top\right]+O(\delta^5).
 \label{eq:compact-root-quartic-main}
\end{align}
In particular, the multiplier error is exactly quintic if \(\chi_2\ne0\).
\end{corollary}

\begin{proof}
Here \(t_0=0\), \(d_R=d_0>0\), and
\[
 \mathcal Z(\varepsilon,0)=Z_\varepsilon(A^\star)
 =\Sc(A^\star)-\frac{\varepsilon^2}{2}J_\star+O(\varepsilon^4),
 \qquad
 \partial_t\mathcal Z(0,0)=U_0CV_0^\top.
\]
It follows that
\(H(\varepsilon,0)=-\chi_\star\varepsilon^2/2+O(\varepsilon^4)\).
The analytic implicit solution therefore satisfies
\[
 t(\varepsilon)=\frac{\chi_\star}{2d_0}\varepsilon^2
 +O(\varepsilon^3).
\]
Since \(\lambda_\varepsilon-\lambda^\star=\varepsilon t(\varepsilon)\),
substitution into the analytic blow-up gives
\eqref{eq:compact-root-cubic-main}--\eqref{eq:compact-root-quadratic-main}.
The range--range and null--null blocks in \(\mathcal N_0\) are Frobenius
orthogonal, which proves \eqref{eq:compact-root-direction-norm-main}.
If \(\chi_\star=0\), the separated positive block has the uniform expansion
\[
 \frac{\sigma}{\sqrt{\sigma^2+\varepsilon^2}}
 =1-\frac{\varepsilon^2}{2\sigma^2}
  +\frac{3\varepsilon^4}{8\sigma^4}+O(\varepsilon^6).
\]
Consequently
\(H(\varepsilon,0)=3\chi_2\varepsilon^4/8+O(\varepsilon^6)\), and the
implicit branch satisfies
\(t(\varepsilon)=-3\chi_2\varepsilon^4/(8d_0)+O(\varepsilon^5)\).
This proves \eqref{eq:compact-root-quintic-main}.  Now
\(s=\varepsilon t(\varepsilon)=O(\varepsilon^5)\).  The quartic term of
the positive block and
\(t(\varepsilon)\partial_t\mathcal Z(0,0)\), with
\(\partial_t\mathcal Z(0,0)=U_0CV_0^\top\), give
\eqref{eq:compact-root-quartic-main}.  Indeed, analyticity yields
\(\mathcal Z(\varepsilon,t(\varepsilon))
=\mathcal Z(\varepsilon,0)+t(\varepsilon)\partial_t\mathcal Z(0,0)
+O(\varepsilon|t(\varepsilon)|+t(\varepsilon)^2)\), so every remaining
term is \(O(\varepsilon^5)\).
\end{proof}

\begin{remark}[An odd-order multiplier hierarchy]
\label{rem:compact-root-hierarchy-main}
The cubic and quintic laws are the first two levels of a complete multiplier
hierarchy.  For an integer \(k\ge1\), let
\[
 b_k=(-1)^k4^{-k}\binom{2k}{k},\qquad
 J_k=U_r\Sigma_r^{-2k}V_r^\top,\qquad
 \chi_k=\langle\Theta,J_k\rangle.
\]
If an integer \(p\ge1\) satisfies
\(a=0\) and \(\chi_1=\cdots=\chi_{p-1}=0\), then the binomial expansion
used above and the analytic implicit function theorem give
\begin{equation}
 \lambda_\delta-\lambda^\star
 =-\frac{b_p\chi_p}{d_0}\,\delta^{2p+1}
  +O(\delta^{2p+2}).
 \label{eq:compact-root-hierarchy-main}
\end{equation}
Thus each further orthogonality cancellation raises the multiplier by two
orders; \(p=1\) is \eqref{eq:compact-root-cubic-main} and \(p=2\) is
\eqref{eq:compact-root-quintic-main}.  This hierarchy is local to the fixed
strict crossing and does not assert that the cancellations are generic.
\end{remark}

At a fixed strict rank-deficient minimizer, these results give a sharp
dichotomy.  Compact-polar stationarity holds exactly when \(a=0\), in which
case smoothing has cubic multiplier bias and quadratic direction bias with
the refinements above.  If \(a\ne0\), compact-polar stationarity fails and the
multiplier bias is necessarily \(\Theta(\delta)\); the direction bias is
\(\Theta(\delta)\) when \(\mathcal M_R\ne0\), but may cancel to second order.

\subsection{A transverse corank-one normal form for boundary and crossover}

For the remaining results assume \(m=n=q\),
\(\operatorname{rank}A^\star=q-1\), and let \(u_0,v_0\) be unit left and
right null vectors.  Put
\begin{equation}
 \alpha=u_0^\top\Theta v_0,
 \qquad \beta=|\alpha|,
 \qquad \alpha\ne0.
 \label{eq:alpha-beta-main}
\end{equation}

\begin{lemma}[Analytic kink normal form]
\label{lem:kink-normal-form}
There are real-analytic functions \(\tau,\ell\) near zero such that
\begin{equation}
 \tau(s)=\alpha s+O(s^2),
 \qquad
 \phi(\lambda^\star+s)=\ell(s)+|\tau(s)|.
 \label{eq:kink-normal-form}
\end{equation}
Moreover,
\begin{equation}
 h_\delta(\lambda^\star+s)
 =\ell_\delta'(s)
 +\frac{\tau(s)\tau'(s)}
       {\sqrt{\tau(s)^2+\delta^2}},
 \qquad
 \ell_\delta'(s)=\ell'(s)+O(\delta^2)
 \label{eq:kink-smooth-derivative}
\end{equation}
uniformly on a fixed small neighborhood, where \(\ell_\delta'(s)\) denotes the
contribution to \(h_\delta\) from the separated nonzero singular-value cluster.
With \(a=\ell'(0)\), the exact
subdifferential is \([a-\beta,a+\beta]\).
\end{lemma}

\begin{proof}
Consider the analytic symmetric dilation
\[
 H(s)=\begin{pmatrix}
 0&A^\star+s\Theta\\
 (A^\star+s\Theta)^\top&0
 \end{pmatrix}.
\]
Its zero eigenspace has basis \((u_0,0),(0,v_0)\), and the compression of
\(H'(0)\) is
\(\bigl(\begin{smallmatrix}0&\alpha\\\alpha&0\end{smallmatrix}\bigr)\).
Rellich perturbation theory \citep[Chapter II]{kato1995perturbation} therefore
gives a signed analytic pair \(\pm\tau(s)\) with
\(\tau'(0)=\alpha\).  The small singular value is \(|\tau(s)|\).  The sum of
the separated positive eigenvalue cluster of the dilation is analytic and
defines \(\ell(s)\), proving \eqref{eq:kink-normal-form}.  The nonzero cluster
stays uniformly separated from zero, so the regular expansion used in
\cref{thm:rate-regular} gives
\(\ell_\delta'=\ell'+O(\delta^2)\).  Differentiating the smoothed small mode
proves \eqref{eq:kink-smooth-derivative}; the two one-sided derivatives at
zero are \(a\mp\beta\).
\end{proof}

\subsection{The corank-one strict specialization and boundary kink}

\begin{corollary}[Linear strict-kink displacement, corank one]
\label{thm:rate-strict}
Under \cref{lem:kink-normal-form}, suppose \(|a|<\beta\).  Then
\begin{equation}
 \frac{\lambda_\delta-\lambda^\star}{\delta}
 \longrightarrow
 t_0=-\frac{a}{\beta\sqrt{\beta^2-a^2}}.
 \label{eq:strict-rate-main}
\end{equation}
The displacement is \(O(\delta)\), and it is \(\Theta(\delta)\) when
\(a\ne0\).
\end{corollary}

\begin{proof}
Here \(C\) has the single positive singular value \(\beta\).
Equation \eqref{eq:eta-strict-rate-main} gives
\(\eta_R=|a|/[\beta\sqrt{\beta^2-a^2}]\), so
\cref{thm:rate-strict-general} reduces to \eqref{eq:strict-rate-main}.
\end{proof}

At the boundary, the small-mode expansion must be grouped with the complete
contact-side derivative; expanding the mode separately loses a linear term.
Here \(\phi'_{\rm con}\) denotes the derivative of the full-rank analytic branch
continued from the stated contact side.
The boundary law is not obtained by sending \(|a|\uparrow\beta\) in the
strict formula: \(\eta_R\to\infty\), so the \(\delta\)-scale balance itself
breaks down and is replaced by the \(\delta^{2/3}\) scale below.

\begin{theorem}[Boundary-kink two-thirds law]
\label{thm:rate-boundary}
Suppose \(|a|=\beta\), and choose the sign of the analytic branch in
\cref{lem:kink-normal-form} so that \(\alpha=\beta>0\).  Let the full-rank
contact side be \(s<0\) when \(a=\beta\) and \(s>0\) when \(a=-\beta\), and
let \(\phi_{\rm con}\) denote the analytic branch continued from that side.
Then \cref{lem:kink-normal-form} gives
\begin{equation}
 \phi'_{\rm con}(\lambda^\star+s)=\kappa s+O(s^2),
 \qquad
 \kappa:=
 \begin{cases}
  \ell''(0)-\tau''(0),&a=\beta,\\
  \ell''(0)+\tau''(0),&a=-\beta.
 \end{cases}
 \label{eq:contact-curvature-main}
\end{equation}
Convexity gives \(\kappa\ge0\); assume only the nondegeneracy
\(\kappa>0\).  Then
\begin{equation}
 \lambda_\delta-\lambda^\star
 =-\operatorname{sign}(a)
 \left(\frac{1}{2\beta\kappa}\right)^{1/3}
 \delta^{2/3}(1+o(1)).
 \label{eq:boundary-rate-main}
\end{equation}
\end{theorem}

\begin{proof}
The expansion in \eqref{eq:contact-curvature-main} is a consequence, not an
extra hypothesis.  On the contact side one has
\(\phi_{\rm con}=\ell-\tau\) when \(a=\beta\) and
\(\phi_{\rm con}=\ell+\tau\) when \(a=-\beta\).  Hence its derivative at
zero is respectively \(a-\beta=0\) or \(a+\beta=0\), and analyticity gives
the displayed expansion with the stated second derivative.

It suffices to treat \(a=\beta\); the other case follows after reversing
\(s\) and the signed analytic branch.  Since
\(\tau(0)=0\), \eqref{eq:kink-smooth-derivative} gives
\(h_\delta(\lambda^\star)=\ell_\delta'(0)=a+O(\delta^2)>0\) for all small
\(\delta\), and \(h_\delta\) is strictly increasing; hence the smooth root
lies on \(s<0\).
Put \(\chi(s)=\operatorname{sign}(\tau(s))\tau'(s)\).  Then
\(\chi(s)=-\beta+O(s)\) and
\(\tau(s)^2=\beta^2s^2(1+O(s))\) on that side.  In the regime
\(\delta/|s|\to0\), group the exact contact derivative before expanding:
\begin{align}
 h_\delta(\lambda^\star+s)
 ={}&\phi'_{\rm con}(\lambda^\star+s)
 +\frac{\delta^2}{2\beta s^2}(1+O(|s|)) \\[-0.2em]
 &+O(\delta^2)+O(\delta^4/s^4).
 \label{eq:boundary-balance-main}
\end{align}
Set \(s=-c\delta^{2/3}\) and divide by \(\delta^{2/3}\).  Every remainder
in \eqref{eq:boundary-balance-main} vanishes, and the limiting equation is
\[
 -\kappa c+\frac{1}{2\beta c^2}=0.
\]
Its positive root is \(c_0=(2\beta\kappa)^{-1/3}\).  Evaluating at fixed
\(c_-<c_0<c_+\) gives opposite signs for all small \(\delta\), so scalar
monotonicity brackets the unique root.  This also proves
\(|s_\delta|\asymp\delta^{2/3}\gg\delta\), validating the expansion.
Letting the brackets approach \(c_0\) proves \eqref{eq:boundary-rate-main}.
\end{proof}

The curvature condition is genuine.  For
\(A(\lambda^\star+s)=\diag(1+s,s)\), one has \(a=\beta=1\), but the contact
branch is constant on \(-1<s<0\), so \(\kappa=0\) and there is no isolated
two-thirds balance.

For the orientation \(a=\beta>0\) used in the crossover theorem below,
this contact curvature is exactly
\(\kappa=\ell''(0)-\tau''(0)=\kappa_0\).

For the boundary member of \cref{thm:open-2by2} with \(b_1=b_2=c=1\) and
\(d=0\), one has \(\lambda^\star=1\), \(a=\beta=1/2\), and
\(\phi'_{\rm con}(1+s)=s/\sqrt{s^2+4}=s/2+O(s^3)\) on the left.  Hence
\begin{equation}
 1-\lambda_\delta\sim2^{1/3}\delta^{2/3}.
 \label{eq:boundary-example-main}
\end{equation}

The exponent is not intrinsic to the certificate slice alone: it also records
how the chosen profile approaches the spectral-ball boundary.

\begin{corollary}[The boundary exponent is a tail law]
\label{cor:profile-tail-rate-main}
Under the hypotheses of \cref{thm:rate-boundary}, let
\(\lambda_\delta^\psi\) be any minimizer of the objective induced by an admissible
profile from \cref{prop:smoothing-family-main}.
\begin{enumerate}
 \item If, for some \(p>1\) and \(\gamma_\psi>0\),
 \begin{equation}
  \psi'(x)=1-\gamma_\psi x^{-p}+o(x^{-p})
  \qquad(x\to\infty),
  \label{eq:profile-polynomial-tail-main}
 \end{equation}
 then
 \begin{equation}
  \lambda_\delta^\psi-\lambda^\star
  =-\sign(a)
  \left(\frac{\gamma_\psi}{\kappa\beta^{p-1}}\right)^{1/(p+1)}
  \delta^{p/(p+1)}(1+o(1)).
  \label{eq:profile-tail-rate-main}
 \end{equation}
 \item If \(\psi'\) first reaches one at a finite threshold
 \(L_\psi>0\) and is identically one thereafter, then
 \begin{equation}
  \lambda_\delta^\psi-\lambda^\star
  =-\sign(a)\frac{L_\psi}{\beta}\,\delta+o(\delta).
  \label{eq:saturated-boundary-rate-main}
 \end{equation}
 For the standard Huber profile in \eqref{eq:two-profile-members-main},
 \(L_\psi=1\) and the sharper expansion is
 \begin{equation}
  \lambda_\delta^{\rm hub}-\lambda^\star
  =-\sign(a)\frac{\delta}{\beta}+O(\delta^2).
  \label{eq:huber-boundary-rate-main}
 \end{equation}
\end{enumerate}
Pseudo-Huber has \(p=2\) and \(\gamma_\psi=1/2\), so
\eqref{eq:profile-tail-rate-main} recovers
\eqref{eq:boundary-rate-main}; finite saturation yields the separate linear
boundary regime.
\end{corollary}

\begin{proof}
It is enough to take \(a=\beta>0\), for which the contact side is \(s<0\).
Along a scale with \(\delta/|s|\to0\), the separated positive cluster differs
from its polar limit by \(o(|s|)\) on the scale obtained below.  Using
\(\tau(s)=\beta s+O(s^2)\) and grouping against the exact contact derivative,
\eqref{eq:profile-polynomial-tail-main} gives
\[
 h_{\delta}^{\psi}(\lambda^\star+s)
 =\kappa s+
 \frac{\gamma_\psi\delta^p}{\beta^{p-1}|s|^p}
 +o\!\left(|s|+\frac{\delta^p}{|s|^p}\right).
\]
Putting \(s=-c\delta^{p/(p+1)}\) and bracketing the monotone root yields
\(\kappa c^{p+1}=\gamma_\psi/\beta^{p-1}\), proving the first claim.

If saturation first occurs at \(L_\psi\), the separated modes are already
saturated for all small \(\delta\).  Monotonicity brackets the root immediately
inside the transition edge \(|\tau(s)|=L_\psi\delta\), and therefore
\(|s|/(L_\psi\delta/\beta)\to1\), proving
\eqref{eq:saturated-boundary-rate-main}.  For standard Huber, inside
\(|\tau(s)|<\delta\) the small mode gives the sharper expansion
\[
 h_\delta^{\rm hub}(\lambda^\star+s)
 =\beta+\frac{\beta^2s}{\delta}
   +O(|s|+s^2/\delta).
\]
The monotone root therefore satisfies
\(s=-\delta/\beta+O(\delta^2)\).  Reflection handles \(a=-\beta\).
\end{proof}

\begin{corollary}[Selected-direction upper bounds]
\label{cor:direction-rates-main}
In fixed dimension and under the corresponding hypotheses of
\cref{thm:rate-regular,thm:direction-strict-general,thm:rate-boundary},
\begin{equation}
 \|\Phi_\delta-\Phi_R^\star\|_F
 =O(\delta^2),\qquad O(\delta),\qquad O(\delta^{2/3}),
 \label{eq:direction-rates-main}
\end{equation}
respectively.  These are upper bounds; exact matrix-direction rates require a
corresponding noncancellation condition.  The strict entry holds for arbitrary
rectangular shape and corank; \cref{thm:direction-strict-general} gives its
first-order constant and the sharper \(O(\delta^2)\) rate when \(a=0\).
\end{corollary}

\begin{proof}
On the separated cluster, the polar factor and smoothed spectral map are
analytic in the matrix argument and \(\delta^2\), so their variation is the
multiplier displacement plus \(O(\delta^2)\); this proves the regular entry.
The strict entry is \eqref{eq:strict-direction-first-main}.  At a boundary
kink, both the singular-vector movement and the coefficient
error are \(O(\delta^{2/3})\) because
\(|s|=O(\delta^{2/3})\) and
\(\delta^2/s^2=O(\delta^{2/3})\).
\end{proof}

\subsection{A local strict-to-boundary crossover}

The fixed-instance constants in
\cref{thm:direction-strict-general,prop:strict-constants-closed-main}
deteriorate as \(\beta-|a|\downarrow0\): \(\eta_R\) diverges and \(d_R\)
vanishes.  The result below is the uniform replacement in precisely that
degenerating direction, currently for a transverse square corank-one
crossing.  Thus the boundary law is not confined to one isolated parameter
value; a moving normal form controls the entire nearby transition.

\begin{lemma}[Moving crossing for general matrix normals]
\label{lem:moving-crossing-main}
Let \(A(\lambda,z)=G(z)+\lambda\Theta(z)\), \(z\in\mathbb R^p\), where
\(G\) and \(\Theta\) are real analytic and \(\Theta(z)\ne0\).  Suppose
the matrices are square, \(A(\lambda_0,0)\) has corank one, and
\begin{equation}
 u_0^\top\Theta(0)v_0\ne0.
 \label{eq:transverse-family-main}
\end{equation}
After shrinking the neighborhood, there are an analytic crossing
\(\lambda_{\rm sing}(z)\) and jointly analytic functions \(\tau,\ell,w\),
with \(s=\lambda-\lambda_{\rm sing}(z)\), such that
\begin{equation}
 \|A(\lambda_{\rm sing}(z)+s,z)\|_*
 =\ell(s,z)+|\tau(s,z)|,
 \qquad
 \tau(s,z)=s\sqrt{w(s,z)},\qquad w(s,z)>0.
 \label{eq:moving-normal-form-main}
\end{equation}
The sign can be chosen so that
\(\beta(z)=\partial_s\tau(0,z)>0\), and \(\ell\) is analytic even if
singular values inside the separated positive cluster cross.
\end{lemma}

\begin{proof}
At a corank-one matrix the adjugate is a nonzero multiple of
\(v_0u_0^\top\).  Condition \eqref{eq:transverse-family-main} makes
\(\partial_\lambda\det A(\lambda_0,0)\ne0\), so the analytic implicit
function theorem gives \(\lambda_{\rm sing}(z)\).  Set
\(\widetilde A(s,z)=A(\lambda_{\rm sing}(z)+s,z)\) and
\(B(s,z)=\widetilde A(s,z)^\top\widetilde A(s,z)\).  For each nearby \(z\),
\(B(0,z)\) has a simple zero eigenvalue separated from the rest; its analytic
continuation \(\mu(s,z)\) satisfies \(\mu(0,z)=0\).  Analytic division gives
\(\mu=s\mu_1\).  Since \(B\succeq0\), the scalar function
\(s\mapsto\mu(s,z)\) has a minimum at zero, so a second division gives
\(\mu=s^2w\).  Transversality makes \(w>0\) after shrinking.

For the remaining modes, a fixed contour around the positive cluster of the
symmetric dilation defines a jointly analytic Riesz projection.  The trace of
the dilation against this projection is the analytic eigenvalue sum
\(\ell(s,z)\), proving \eqref{eq:moving-normal-form-main}.
\end{proof}

\begin{theorem}[Unified cubic-root crossover]
\label{thm:crossover-main}
In \cref{lem:moving-crossing-main}, suppose at \(z=0\)
\begin{equation}
 a_0=\partial_s\ell(0,0)=\beta_0>0,
 \qquad
 \kappa_0=\partial_{ss}\ell(0,0)-\partial_{ss}\tau(0,0)>0.
 \label{eq:crossover-boundary-data-main}
\end{equation}
Let \(z=z_\delta\to0\) approach from the strict or boundary side and set
\begin{align*}
 a_\delta&=\partial_s\ell(0,z_\delta),
 &\beta_\delta&=\partial_s\tau(0,z_\delta),\\
 \Delta_\delta&=\beta_\delta-a_\delta\ge0,
 &\kappa_\delta&=\partial_{ss}\ell(0,z_\delta)
                    -\partial_{ss}\tau(0,z_\delta).
\end{align*}
Let \(\lambda_\delta^\star=\lambda_{\rm sing}(z_\delta)\), let
\(\lambda_\delta\) be the smoothed minimizer, and put
\(x_\delta=\lambda_\delta^\star-\lambda_\delta\).  Define
\(\widehat x_\delta>0\) as the unique positive root of
\begin{equation}
 \kappa_\delta\widehat x_\delta^3
 +\Delta_\delta\widehat x_\delta^2
 =\frac{\delta^2}{2\beta_\delta}.
 \label{eq:cubic-root-main}
\end{equation}
Then \(x_\delta>0\) for all small \(\delta\), and, along every such path,
including oscillatory ones,
\begin{equation}
 \frac{x_\delta}{\widehat x_\delta}
 =1+O(\Delta_\delta+\widehat x_\delta).
 \label{eq:cubic-relative-main}
\end{equation}
Consequently:
\begin{enumerate}
 \item if \(\Delta_\delta/\delta^{2/3}\to\zeta\in[0,\infty)\), then
 \begin{equation}
  \frac{x_\delta}{\delta^{2/3}}\to y_\zeta,
  \qquad
  \kappa_0y_\zeta^3+\zeta y_\zeta^2=\frac{1}{2\beta_0};
  \label{eq:crossover-critical-main}
 \end{equation}
 \item if \(\Delta_\delta/\delta^{2/3}\to\infty\), then
 \begin{equation}
  \frac{x_\delta\sqrt{\Delta_\delta}}{\delta}
  \to\frac{1}{\sqrt{2\beta_0}};
  \label{eq:crossover-inner-main}
 \end{equation}
 \item for arbitrary paths,
 \begin{equation}
  x_\delta\asymp\widehat x_\delta\asymp
  \min\left\{\frac{\delta}{\sqrt{\Delta_\delta}},
              \delta^{2/3}\right\},
  \label{eq:crossover-scale-main}
 \end{equation}
 where the first entry is \(+\infty\) at \(\Delta_\delta=0\).
\end{enumerate}
\end{theorem}

\begin{proof}
The one-sided derivatives at the moving crossing are
\(-\Delta_\delta\le0\) and
\(a_\delta+\beta_\delta\to2\beta_0>0\), so
\(\lambda_\delta^\star\) is an exact minimizer.  For the instance indexed by
\(z\), set
\[
 h_{\delta,z}(\lambda)
 =\langle\Theta(z),Z_\delta(A(\lambda,z))\rangle,
 \quad
 \Delta(z)=\beta(z)-a(z),
 \quad
 \kappa(z)=\partial_{ss}\ell(0,z)-\partial_{ss}\tau(0,z),
\]
and \(H_{\delta,z}(x)=h_{\delta,z}(\lambda_{\rm sing}(z)-x)\).  Joint
analyticity and the grouped expansion in \eqref{eq:boundary-balance-main}
give constants \(z_0,x_0,\eta_0,C>0\), independent of the path, such that
whenever \(\|z\|\le z_0\), \(0<x\le x_0\), and
\(0<\delta\le\eta_0x\),
\begin{align}
 H_{\delta,z}(x)
 ={}&-\Delta(z)-\kappa(z)x
 +\frac{\delta^2}{2\beta(z)x^2}+E_{\delta,z}(x),
 \label{eq:crossover-functional-main}\\[-0.2em]
 |E_{\delta,z}(x)|
 \le{}&C\left(x^2+\frac{\delta^2}{x}
                    +\delta^2+\frac{\delta^4}{x^4}\right).
 \label{eq:crossover-remainder-main}
\end{align}
At \(x=0\), \(H_{\delta,z_\delta}(0)=a_\delta+O(\delta^2)>0\), whereas
it is negative at a sufficiently small fixed \(x_0>0\).  Strict convexity
places the unique root in \((0,x_0)\).

Put
\(M_\delta=\Delta_\delta+\kappa_\delta\widehat x_\delta\).  Equation
\eqref{eq:cubic-root-main} is equivalent to
\[
 \frac{\delta^2}{2\beta_\delta\widehat x_\delta^2}=M_\delta.
\]
It implies \(\widehat x_\delta,M_\delta\to0\) and
\(\delta/\widehat x_\delta=\sqrt{2\beta_\delta M_\delta}\to0\), so the
uniform expansion applies at \(x=c\widehat x_\delta\).  Uniformly for \(c\)
in a compact subset of \((0,\infty)\), the four terms in the remainder,
divided by \(M_\delta\), are respectively
\[
 O(\widehat x_\delta),\quad O(\widehat x_\delta),\quad
 O(\widehat x_\delta^2),\quad O(M_\delta).
\]
The leading part at this point is
\[
 \Delta_\delta(c^{-2}-1)
 +\kappa_\delta\widehat x_\delta(c^{-2}-c).
\]
It is a fixed positive multiple of \(M_\delta\) for fixed \(c<1\), and a
fixed negative multiple for fixed \(c>1\).  Monotonicity therefore gives
\(x_\delta/\widehat x_\delta\to1\).  More quantitatively, for
\(c=1+u\) with \(|u|\) small, the leading term has the sign of \(-u\) and
magnitude at least \(c_0|u|M_\delta\), whereas the remainder is bounded by
\(C_1(\Delta_\delta+\widehat x_\delta)M_\delta\), with uniform positive
constants \(c_0,C_1\).  Taking
\(u=\pm K(\Delta_\delta+\widehat x_\delta)\) and
\(K>C_1/c_0\) gives the quantitative bracket
\eqref{eq:cubic-relative-main}.

The three consequences now follow directly from
\eqref{eq:cubic-root-main}: divide it by \(\delta^2\) in the critical window,
drop its cubic term when \(\Delta_\delta/\delta^{2/3}\to\infty\), and compare
the two nonnegative terms for the two-sided order bound.
\end{proof}

The cubic also controls the matrix actually returned by the oracle.  The
comparator must be the center of the moving instance rather than the fixed
center at \(z=0\).

\begin{corollary}[Direction-level crossover]
\label{cor:direction-crossover-main}
Under \cref{thm:crossover-main}, let \(P(s,z)\) be the analytic polar
contribution of the separated positive singular cluster and let \(D(s,z)\) be
the analytic signed dyad of the small mode, chosen so that
\(\langle\Theta(z),D(0,z)\rangle=\beta(z)>0\).  With derivatives evaluated at
\((s,z)=(0,0)\), define
 \begin{equation}
 \mathcal M_0=-\partial_sP(0,0)+\partial_sD(0,0)
 +\frac{\kappa_0}{\beta_0}D(0,0).
 \label{eq:direction-M0-main}
\end{equation}
Then, relative to the exact Fenchel center \(\Phi_R^\star(z_\delta)\) of the
moving instance,
\begin{equation}
 \frac{\Phi_\delta-\Phi_R^\star(z_\delta)}{x_\delta}
 \longrightarrow\mathcal M_0.
 \label{eq:direction-crossover-limit-main}
\end{equation}
Hence
\[
 \|\Phi_\delta-\Phi_R^\star(z_\delta)\|_F
 =O(\widehat x_\delta),
\]
and if \(\mathcal M_0\ne0\), this norm is asymptotic to
\(\|\mathcal M_0\|_F\widehat x_\delta\); if \(\mathcal M_0=0\), it is
\(o(\widehat x_\delta)\).
\end{corollary}

\begin{proof}
The smoothed spectral decomposition at \(s=-x_\delta\) is
\[
 \Phi_\delta=P_\delta(-x_\delta,z_\delta)
 +r_\delta(-x_\delta,z_\delta)D(-x_\delta,z_\delta),
 \quad
 r_\delta=\frac{\tau}{\sqrt{\tau^2+\delta^2}},
\]
where \(P_\delta=P+O(\delta^2)\).  Tangency gives
\(r_\delta=-\ell_\delta'/\tau'\).  Expanding this quotient and using
\(a=\beta-\Delta\) and
\(\kappa=\ell_{ss}-\tau_{ss}\) gives, uniformly,
\[
 r_\delta(-x,z)+\frac{a(z)}{\beta(z)}
 =x\left\{\frac{\kappa(z)}{\beta(z)}
 +\frac{\Delta(z)\tau_{ss}(0,z)}{\beta(z)^2}\right\}
 +O(x^2+\delta^2).
\]
Since \(\Delta(z_\delta)\to0\) and \(\delta^2/x_\delta\to0\), this yields
\[
 \frac{r_\delta(-x_\delta,z_\delta)+a(z_\delta)/\beta(z_\delta)}
 {x_\delta}\longrightarrow\frac{\kappa_0}{\beta_0}.
\]
Since
\(\Phi_R^\star(z)=P(0,z)-[a(z)/\beta(z)]D(0,z)\), Taylor expansion proves
\eqref{eq:direction-crossover-limit-main}; the separated-cluster smoothing
error is negligible because \(\delta^2/x_\delta\to0\).  The norm statements
follow from \eqref{eq:cubic-relative-main}.
\end{proof}

The moving-crossing recentering is essential: nonnegativity together with a
single identity \(\mu(0,0)=0\) would not imply the factorization
\(\mu=s^2w\).  The theorem is uniform only near one fixed analytic transverse
corank-one boundary instance.

\section{Numerical verification of the constants}
\label{sec:numerics}

The local laws predict constants, not only exponents.  Assertion-based scripts
therefore test closed-form instances at high precision; the full tables and
execution details are in the technical supplement.  For the fixed
\(2\times2\) family
\(G=\bigl(\begin{smallmatrix}0&b_1\\b_2&c\end{smallmatrix}\bigr)\),
\(\Theta=e_1e_1^\top\), the three basic normalized limits are:
\begin{center}
\begin{tabular}{@{}lll@{}}
\toprule
Regime & predicted & computed at the smallest \(\delta\)\\
\midrule
regular & \(-16/9=-1.777777778\) & \(-1.777777777\)\\
strict corank one & \(-5/(4\sqrt{15})=-0.3227486122\) & \(-0.3227486102\)\\
boundary & \(2^{1/3}=1.259921050\) & \(1.259921042\)\\
\bottomrule
\end{tabular}
\end{center}
On the same boundary instance, standard Huber instead gives
\((1-\lambda_\delta^{\rm hub})/\delta=1.99999997\) at
\(\delta=10^{-8}\), versus the prediction \(1/\beta=2\) in
\eqref{eq:huber-boundary-rate-main}.
For the admissible profile
\(\psi'(x)=x/(1+x^3)^{1/3}\) on \(x\ge0\), the predicted \(p=3\)
constant is \(1.2778862085\); the computed normalized displacement is
\(1.2778860473\) at \(\delta=10^{-9}\).  This checks that the exponent
changes with the tail, not merely between the two named profiles.

For a corank-two check, take
\(A^\star=\diag(1,1,0,0)\),
\(\Theta=\diag(3/8,3/8,1,1/2)\), and \(\lambda^\star=1\).
Here \(\eta_R=0.715080068980\) and
\(\|K_F-K_R\|_F=0.040996138377\).  At \(\delta=10^{-4}\), the
pseudo-Huber slope is \(-0.7150801\), while Huber gives the water-filling
slope \(-0.6000000000\) and recovers \(\Phi_F^\star\) to \(10^{-11}\).
Because \(P_{\rm big}(s)\) and \(\widetilde B(s)\) are constant on this
diagonal line, \(\partial_s\mathcal W=0\) proves
\(t_1=\mathcal M_R=0\).  A non-diagonal rational instance instead predicts
\(\|\mathcal M_R\|_F=0.671775381\), versus the scaled value
\(0.671776068\) at \(\delta=10^{-6}\).

The compact-root instance
\(A^\star=\diag(2,1,0)\), \(\Theta=\diag(1,-1,1)\) verifies the predicted
cubic coefficient \(-3/8\) and quadratic direction norm \(\sqrt{26}/8\).
After the second cancellation, the rational \(4\times4\) instance in the
supplement has \(\chi_2=1/3\) and converges monotonically to the quintic
coefficient \(-1/8\); its normalized quartic direction converges to the full
matrix in \eqref{eq:compact-root-quartic-main}.

The moving-family assertions cover boundary, critical, strict-side, and
oscillatory paths and verify both the cubic root and direction constants.
Complete data, matrices, and finite-scale sequences are in the supplement;
the checks test constants but imply no occurrence-frequency law.

\section{Discussion and limitations}
\label{sec:discussion}

The two remedies make different tradeoffs.  Exact interval bisection followed
by \eqref{eq:KF-main} has no smoothing bias but needs a rank decision.
Pseudo-Huber gives a strictly monotone equation and converges to \(K_R\),
whereas the family in \cref{prop:smoothing-family-main} identifies \(K_F\) as
the Huber Fenchel center.  These centers agree only in the cases of
\cref{cor:center-agreement-main}.  The boundary bias is profile dependent:
polynomial tail order \(p\) gives \(\delta^{p/(p+1)}\), pseudo-Huber gives
\(\delta^{2/3}\), and finite saturation gives \(\delta\).  At fixed strict
crossings the first direction coefficient is computable from
\eqref{eq:frozen-field-main}; compact-root instances instead begin at
cubic/quadratic order and obey the multiplier hierarchy of
\cref{rem:compact-root-hierarchy-main}.

The strict laws allow arbitrary shape and corank, but boundary and crossover
laws currently require a transverse square corank-one crossing; multiple
signed small modes and nontransverse exponents remain open.  The admissible
family does not cover every vanishing regularizer, and full-space openness
implies no frequency law without a distribution for \((G,\Theta)\).  No outer
complexity or learning-performance gain follows without a separate descent
model and control of the changing normal.

\section{Conclusion}

We characterized the tangency-constrained certificate slice for nuclear-norm
approximation from a matrix line and identified its hard and Fenchel spectral
centers.  An admissible smoothing family links center selection and boundary
exponents, while strict arbitrary-corank paths admit computable analytic
multiplier and direction expansions.  Transverse corank-one boundaries
exhibit the sharp pseudo-Huber two-thirds law and a uniform cubic crossover.
Compact-polar open failure explains why the full certificate slice is
essential.

\section*{Data and code availability}

No external data are used.  Assertion-based Python scripts reproducing every
numerical table and constant are included in the supplementary source archive
submitted with the manuscript; the archive also records the required Python,
NumPy, SciPy, and mpmath dependencies.

\section*{Declaration on generative AI assistance}

Generative AI tools assisted with language editing, literature discovery,
algebraic exploration, drafting portions of theorem and proof text, and
drafting verification code.  The authors independently checked every cited
source, proof, and computational result and assume responsibility for all
content.


{\footnotesize
\setlength{\parskip}{0pt}
\setlength{\bibsep}{0pt plus 0.2ex}
\bibliographystyle{plainnat}
\bibliography{references}
}

\end{document}
